% concatenated from Morse_theory_and_moduli_spaces_of_self_avoiding_polygonal_linkages.tex
 %-----------------------------------------------------------------------
% Beginning of tran-l-template.tex
%-----------------------------------------------------------------------
%
%     This is a topmatter template file for TRAN for use with AMS-LaTeX.
%
%     Templates for various common text, math and figure elements are
%     given following the \end{document} line.
%
%%%%%%%%%%%%%%%%%%%%%%%%%%%%%%%%%%%%%%%%%%%%%%%%%%%%%%%%%%%%%%%%%%%%%%%%

%     Remove any commented or uncommented macros you do not use.

  %\documentclass{amsart}
\documentclass[11pt,letterpaper]{amsart}

%\documentclass[11pt,letterpaper,draft]{amsart}
%     If you need symbols beyond the basic set, uncomment this command.
%\usepackage{amssymb}

%     If your article includes graphics, uncomment this command.
%\usepackage{graphicx}

%     If the article includes commutative diagrams,
%\usepackage[cmtip,all]{xy}


%     Update the information and uncomment if AMS is not the copyright
%     holder.
%\copyrightinfo{2009}{American Mathematical Society}

%%%%%%%%%%%%%%%%%%%%%
%%%鑷娣诲姞%%Begin

 \usepackage{txfonts} %%%%%%%%瀛椾綋 $\lambda$
\usepackage{graphicx}
\usepackage{verbatim}
 %\usepackage{color}
\usepackage{tikz}
  %\usepackage[switch]{lineno} \linenumbers
 % \usepackage{lineno} \linenumbers

\usepackage{caption}
%
\usepackage{subcaption}


\usepackage[colorlinks,linkcolor=blue,citecolor=blue]{hyperref}

 \theoremstyle{plain}

  %\usepackage{bibentry}


%%%%%%%%%%%%%%%%%%%%%
%%%鑷娣诲姞%%End


%%%%%%%%%%%%%%%%%%%%%

\newtheorem{theorem}{Theorem}[section]
\newtheorem{lemma}[theorem]{Lemma}
\newtheorem*{claim}{Claim}
\newtheorem{problem}[theorem]{Problem}

\newtheorem{corollary}[theorem]{Corollary}
\newtheorem{proposition}[theorem]{Proposition}

%\newtheorem{remarks}[theorem]{Remarks}
\newtheorem{conjecture}[theorem]{Conjecture}

\theoremstyle{definition}
\newtheorem{definition}[theorem]{Definition}
\newtheorem{remark}[theorem]{Remark}
%\newtheorem{example}[theorem]{Example}

%\newtheorem{xca}[theorem]{Exercise}


 %%%%%%%%%%%%%%%%%%%%%%%%%%
 \newcommand{\ba}{\bar{a}}
 \newcommand{\bb}{\bar{b}}
 \newcommand{\bi}{\bar{i}}
 \newcommand{\bj}{\bar{j}}
 \newcommand{\bk}{\bar{k}}
 \newcommand{\bl}{\bar{l}}
 \newcommand{\bm}{\bar{m}}
 \newcommand{\bn}{\bar{n}}
 \newcommand{\bp}{\bar{p}}
 \newcommand{\bq}{\bar{q}}
 \newcommand{\br}{\bar{r}}
 \newcommand{\bs}{\bar{s}}
 \newcommand{\bt}{\bar{t}}
 \newcommand{\bw}{\bar{w}}
 \newcommand{\bz}{\bar{z}}

 %\newcommand{\v1}{{\bf 1}}
 \newcommand{\vn}{{\bf n}}

 \newcommand{\bM}{\bar{M}}


 \newcommand{\balpha}{\bar{\alpha}}
 \newcommand{\bbeta}{\bar{\beta}}
 \newcommand{\eeta}{\bar{\eta}}
 \newcommand{\btau}{\bar{\tau}}
 \newcommand{\bxi}{\bar{\xi}}
 \newcommand{\bzeta}{\bar{\zeta}}

 \newcommand{\bpartial}{\bar{\partial}}

 %%%%%%%%%%%%%%%%%%%%%%%%%%

 \newcommand{\cA}{\mathcal{A}}
 \newcommand{\cB}{\mathcal{B}}
 \newcommand{\cF}{\mathcal{F}}
 \newcommand{\cH}{\mathcal{H}}
 \newcommand{\cK}{\mathcal{K}}
 \newcommand{\cL}{\mathcal{L}}
 \newcommand{\cM}{\mathcal{M}}
 \newcommand{\cT}{\mathcal{T}}
 \newcommand{\Hess}{\mbox{Hess}}
 \newcommand{\Ric}{\mbox{Ric}}
 \newcommand{\tr}{\mbox{tr}}

 \newcommand{\tF}{\tilde{F}}
 \newcommand{\tU}{\tilde{U}}

 %%%%%%%%%%%%%%%%%%%%%%%%%%%%%%%%%%%%%%%


\theoremstyle{remark}
%\newtheorem{remark}[theorem]{Remark}
\newtheorem*{example}{Example}%[theorem][theorem]
\numberwithin{equation}{section}


\makeatletter
\@namedef{subjclassname@1991}{2020 Mathematics Subject Classification}
\makeatother
\begin{document}

\captionsetup[figure]{labelfont={bf},labelformat={default},labelsep=period,name={Fig.}}


\title[Moduli spaces of self-avoiding linkages]{Morse theory and moduli spaces of self-avoiding polygonal linkages}


\author{Te Ba}
\address{School of Mathematics, Hunan University, Changsha 410082, China}
\email{batexu@hnu.edu.cn}
\author{Ze Zhou}
\address{School of Mathematical Sciences, Shenzhen University,  Shenzhen 518060, China }
\email{zhouze@szu.edu.cn}
%\thanks{The author is supported by NSFC (No.12371075 and No.11631010).}
%\textsc{Ze Zhou} \\
%\normalsize School of Mathematical Sciences, Shenzhen University,  Shenzhen 518060, China \\
%\normalsize{zhouze@szu.edu.cn}
%}

\date{}

\subjclass{Primary 52C25; Secondary 57R70, 68U05.}%55U05,
%, 68T40 70B15, 55R80


\begin{abstract}
We show that a smooth $d$-manifold $M$ is diffeomorphic to $\mathbb R^d$ if it admits a Lyapunov-Reeb function, i.e., a smooth map $f:M\to\mathbb R$ that is proper, lower-bounded, and has a unique critical point. By constructing such functions, we prove that the moduli spaces of self-avoiding polygonal linkages and configurations are diffeomorphic to Euclidean spaces. This resolves the Refined Carpenter's Rule Problem and confirms a conjecture proposed by Gonz\'{a}lez and Sedano-Mendoza. Furthermore, we describe foliation structures of these moduli spaces via level sets of Lyapunov-Reeb functions and develop algorithms for related problems.
\end{abstract}

\maketitle
%\tableofcontents

\section{Introduction}\label{S-1}
Let $G(V,E)$ be a simple planar graph and let $\ell:E\to\mathbb{R}_+$ be an assignment of positive reals to the edges of $G$. A \textbf{configuration} for $G$ in $\mathbb R^2$ is a map $p:G\to\mathbb R^{2}$ such that the image $p(e)$ is a line segment for each $e\in E$, and a \textbf{linkage} for $(G,\ell)$ in $\mathbb R^2$ is a configuration $p: G\to \mathbb R^{2}$ such that $p(e)$ has length $\ell(e)$ for every $e\in E$.  Over the past decades, linkage theory has emerged as a foundational framework across diverse disciplines, including discrete and computational geometry~\cite{Hopcroft-1985,Alt-2003}, topological robotics~\cite{Lumelsky-1987}, molecular biology~\cite{Aichholzer-2003}, polymer physics~\cite{Frank-Kamenetskii-1997}, and related fields. For additional background, we refer readers to the seminal monographs by Demaine and O'Rourke~\cite{Demaine-2007}, Farber~\cite{Farber-2008}, and Connelly and Demaine~\cite{Connelly-2009}.

We say a configuration $p: G\to \mathbb R^{2}$ is \textbf{self-avoiding} if it is an embedding. The moduli space $\mathcal M(G)$ is defined to be the set of equivalence classes of all self-avoiding configurations for $G$ in $\mathbb R^2$, where two configurations $p_1,p_2: G\to \mathbb R^{2}$ are equivalent if there exists an orientation-preserving isometry
$\zeta:\mathbb R^{2}\to \mathbb R^{2}$ such that $p_2=\zeta\circ p_1$.  For $p_1,p_2\in\mathcal M(G)$, the distance between them is defined as
\[
d(p_1,p_2)=\inf\nolimits_{\zeta}\left(\sum\nolimits_{v\in V}\|p_2(v)-\zeta\circ p_1(v)\|\right),
\]
where the  infimum is taken over all orientation-preserving isometries of $\mathbb R^2$. Similarly, we define $\mathcal M(G,\ell)\subset \mathcal M(G)$ as the subset of equivalence classes of self-avoiding linkages for $(G,\ell)$ in $\mathbb R^2$ and endow it with the induced metric.

In this paper we restrict our attention to polygonal linkages and configurations. Specifically, we assume $G=R_m$ with $m\geq2$ or $G=C_m$ with $m\geq3$, where $R_m$ (resp. $C_m$) denotes the graph with vertices $v_1,\cdots, v_m$ such that an edge exists between $v_i$ and $v_j$ if and only if
\[
|j-i|=1\quad \left(\text{resp}.\;\;|j-i|=1\;\;\text{or}\;\;|j-i|=m-1\right).
\]
Up to orientation-preserving isometries, the moduli space $\mathcal M(R_m)$
(resp. $\mathcal M(C_m)$) parameterizes all simple polylines (resp. polygons) in
$\mathbb R^2$ with $m$ marked vertices. Accordingly, the subspace
$\mathcal M(R_m,\ell)\subset \mathcal M(R_m)$
(resp. $\mathcal M(C_m,\ell)\subset \mathcal M(C_m)$) parameterizes simple polylines (resp. polygons) with side lengths prescribed by $\ell$. We call an element in $\mathcal M(R_m)$ (resp. $\mathcal M(R_m,\ell)$) an \textbf{arm configuration} (resp. \textbf{arm linkage}) and call an element in $\mathcal M(C_m)$ (resp. $\mathcal M(C_m,\ell)$) a \textbf{cycle configuration} (resp. \textbf{cycle linkage}), as illustrated in Fig.~\ref{F-1}.


\begin{figure}[htbp]
     \centering
     \begin{minipage}[b]{0.40\textwidth}
     \centering
\begin{tikzpicture}[scale=0.66]
\coordinate (a) at (-0.5,-0.4);
\coordinate (b) at (1.7,-1);
\coordinate (c) at (1,-3);
\coordinate (d) at (2.5,-3.4);
\coordinate (e) at (1,-4.5);
\coordinate (f) at (-1,-4);
\coordinate (g) at (-2,-2.18);
\coordinate (h) at (-1,-1);
\coordinate (i) at (0,-1.4);
\draw (a)--(b)--(c)--(d)--(e)--(f)--(g)--(h)--(i);%--(j)
\fill (a) circle (1.5pt);
\fill (b) circle (1.5pt);
\fill (c) circle (1.5pt);
\fill (d) circle (1.5pt);
\fill (e) circle (1.5pt);
\fill (f) circle (1.5pt);
\fill (g) circle (1.5pt);
\fill (h) circle (1.5pt);
\fill (i) circle (1.5pt);
\end{tikzpicture}
     \end{minipage}
     \hfill
     \begin{minipage}[b]{0.50\textwidth}
     \centering
\begin{tikzpicture}[scale=0.75]
\coordinate (a) at (-0.7,0.2);
\coordinate (b) at (1,-1);
\coordinate (c) at (2.55,0.63);
\coordinate (d) at (3,-2.4);
\coordinate (e) at (2,-1.4);
\coordinate (f) at (1.38,-3);
\coordinate (g) at (0.38,-1.7);
\coordinate (h) at (-1.32,-2.18);
\coordinate (i) at (-0.62,-1.18);

\draw (a)--(b)--(c)--(d)--(e)--(f)--(g)--(h)--(i)--(a);
\fill (a) circle (1.5pt);
\fill (b) circle (1.5pt);
\fill (c) circle (1.5pt);
\fill (d) circle (1.5pt);
\fill (e) circle (1.5pt);
\fill (f) circle (1.5pt);
\fill (g) circle (1.5pt);
\fill (h) circle (1.5pt);
\fill (i) circle (1.5pt);
\end{tikzpicture}
     \end{minipage}
     \hfill
\caption{The left is an arm linkage and the right is a cycle linkage}\label{F-1}
\end{figure}

A fundamental problem in linkage theory is to investigate the topology of moduli spaces. For instance, the connectivity of $\mathcal M(R_m,\ell)$ or $\mathcal M(C_m,\ell)$ is closely tied to the Carpenter's Rule Problem: \emph{Is it always possible to deform any simple polyline to a line segment} (by G. Bergman), \emph{or to deform any simple polygon to a convex polygon} (by S. Schanuel), \emph{such that no self-intersection is created and each side remains a fixed length during the deformation process?} It was mentioned in Kirby's list of \textsl{Problems in Low-Dimensional Topology} (see~\cite[Problem 5.18]{Kirby-1997}) and remained open for a long time until Connelly, Demaine and Rote~\cite{Connelly-2003} provided a celebrated answer in 2003. An alternative approach, along with an efficient algorithm, was developed by Streinu~\cite{Streinu-2005}. Additionally, Lenhart and Whitesides~\cite{Lenhart-1995} and Aichholzer et al.~\cite{Aichholzer-2001} described the topology of subspaces of convex cycle linkages. The results combined lead to the following theorem.

\begin{theorem}\label{T-1-1} Let $R_m$ and $C_m$ be defined as above. The following hold:
\begin{itemize}
\item[$(i)$] $\mathcal M(R_m,\ell)$ is path-connected.
\item[$(ii)$] $\mathcal M(C_m,\ell)$ is either empty or consists of exactly two connected components, where each linkage in one component is the reflection of a linkage in the other component.
\end{itemize}
\end{theorem}

\begin{remark}\label{R-1-2}
By a theorem of Penner~\cite[Theorem 6.2]{Penner-1987}, $\mathcal M(C_m,\ell)\neq \emptyset$ if and only if $\ell:E\to \mathbb R_{+}$ satisfies the following condition:
\begin{itemize}
\item[$\mathrm{\mathbf{(c1)}}$]For every $e_i\in E$, $\ell(e_i)<\sum_{e_j\in E\setminus\{e_i\}} \ell(e_j)$.
\end{itemize}
\end{remark}

It is of interest to explore topological properties of these moduli spaces in greater depth. Let $\mathcal M^{+}(C_m,\ell)\subset \mathcal M(C_m,\ell)$ be the component of positively oriented cycle linkages. Connelly, Demaine and Rote~\cite{Connelly-2003} claimed $\mathcal M(R_m,\ell)$ and $\mathcal M^{+}(C_m,\ell)$ are contractible, with proofs later provided in Farber's monograph~\cite{Farber-2008} and research of Han et al.~\cite{Han-2013} and Shimamoto and Wootters~\cite{Shimamoto-2014}. Strikingly, Shimamoto and Wootters~\cite{Shimamoto-2014} further deduced that $\mathcal M^{+}(C_m,\ell)$ is homeomorphic to $\mathbb R^{m-3}$
under condition~$\mathrm{\mathbf{(c1)}}$ and the constraint that $\sum_{e_i\in E}\sigma(e_i)\ell(e_i)\neq0$ for every sign assignment $\sigma: E\to\{-1,1\}$ to the edges.


These results motivate the Refined Carpenter's Rule Problem for cycle linkages: \emph{Is $\mathcal M^{+}(C_m,\ell)$ homeomorphic to $\mathbb R^{m-3}$ whenever $\mathcal M^{+}(C_m,\ell)\neq \emptyset$?} This problem is tantamount to asking whether the constraint stipulated in Shimamoto and Wootters' result can be removed. The core challenge arises when a sequence of linkages in $\mathcal M^{+}(C_m,\ell)$ approaches a pair of overlapping line segments (see Remark~\ref{R-4-13}).
An analogous question can also be posed for the moduli space $\mathcal M(R_m,\ell)$.

Recently, characterizing the topology of the moduli space $\mathcal M(C_m)$ has garnered significant attention. In~\cite{Gonzalez-2018}, Gonz\'{a}lez proved that $\mathcal M(C_m)$ consists of two homeomorphic components, each of which is simply connected. Let $\mathcal M^{+}(C_m)\subset \mathcal M (C_m)$ be the component corresponding to positively oriented polygons. Subsequently, Gonz\'{a}lez and Sedano-Mendoza~\cite[Conjecture 1.1]{Gonzalez-2023} proposed the following conjecture: \emph{$\mathcal M^{+}(C_m)$ is homeomorphic to $\mathbb R^{2m-3}$}. For $m=4$, this
conjecture was validated by Gonz\'{a}lez and  L\'{o}pez-L\'{o}pez~\cite{Gonzalez-2016}. Thereafter, Fortier Bourque~\cite{Fortier Bourque-2024} investigated the moduli spaces of marked immersed polygons and solved the problem for $m=5$. To date, the general case for $m>5$ remains open.

This paper aims to address the aforementioned problems. We first note that the moduli spaces $\mathcal M(R_m)$, $\mathcal M(R_m,\ell)$, $\mathcal M^{+}(C_m)$, and $\mathcal M^{+}(C_m,\ell)$ are smooth (i.e., $C^\infty$) manifolds (see Lemma~\ref{L-3-1} and Lemma~\ref{L-4-1}). Leveraging smooth structures, we further give a complete description of the topology of these moduli spaces. In what follows, the Euclidean $d$-space $\mathbb R^d$ and the $d$-sphere $S^d$ are assumed to carry their standard smooth structures.


\begin{theorem}\label{T-1-3}
Suppose $m\geq 2$ and suppose $\ell:E\to\mathbb R_{+}$ is an arbitrary assignment. We have the following conclusions:
\begin{itemize}
\item[$(i)$] $\mathcal M(R_m,\ell)$ is diffeomorphic to $\mathbb{R}^{m-2}$.
\item[$(ii)$]$\mathcal M(R_m)$ is diffeomorphic to $\mathbb{R}^{2m-3}$.
\end{itemize}
\end{theorem}


\begin{theorem}\label{T-1-4}%\rm{\textbf{(c1)}}
Suppose $m\geq 3$ and suppose $\ell:E\to\mathbb R_{+}$ is an assignment satisfying condition $\mathrm{\mathbf{(c1)}}$
(in Remark~\ref{R-1-2}). The following statements hold:
\begin{itemize}
\item[$(i)$]  $\mathcal M^{+}(C_m,\ell)$ is diffeomorphic to $\mathbb{R}^{m-3}$.
\item[$(ii)$] $\mathcal M^{+}(C_m)$ is diffeomorphic to $\mathbb{R}^{2m-3}$.
\end{itemize}
\end{theorem}

\begin{remark}
Theorems~\ref{T-1-3} and~\ref{T-1-4} resolve the Refined Carpenter's Rule Problem, offering a unified generalization of the results of Connelly, Demaine and Rote~\cite{Connelly-2003} and Shimamoto and Wootters~\cite{Shimamoto-2014}. Moreover, part $(ii)$ of Theorem~\ref{T-1-4} confirms the conjecture of Gonz\'{a}lez and Sedano-Mendoza~\cite[Conjecture 1.1]{Gonzalez-2023}.
\end{remark}


For the proofs of the above two results, we shall make use of Morse theory, which has motivated extensive research on linkages allowing self-intersections. Relevant works include those by  Hausmann~\cite{Hausmann-1989}, Kapovich and Millson~\cite{Kapovich-1995}, Kamiyama~\cite{Kamiyama-1998}, Milgram and Trinkle~\cite{Milgram-2004}, Shimamoto and Vanderwaart~\cite{Shimamoto-2005}, Farber and Sch\"{u}tz~\cite{Farber-2007}, Khimshiashvilli~\cite{Khimshiashvili-2009}, Sch\"{u}tz~\cite{Schuetz-2010,Schuetz-2013,Schuetz-2016}, Farber, Hausmann and Sch\"{u}tz~\cite{Farber-Hausmann-2011}, Farber and Fromm~\cite{Farber-Fromm-2012}, among others.

However, to adapt to the study of self-avoiding linkages and configurations, we need to go a step further to develop Morse theory for open manifolds. To start with, we assume $M$ is a smooth manifold.

\begin{definition}\label{D-1-6}
A smooth map $f:M\to \mathbb R$ is called a \textbf{Lyapunov-Reeb function} (or \textbf{L-R function} for short) on $M$ if it is proper, lower-bounded, and has exactly one critical point.
\end{definition}

Recall that a \textbf{proper map} between two topological spaces means the preimage of every compact subset is compact. We establish the following result, analogous to the Reeb Sphere Theorem~\cite{Reeb-1946,Milnor-1964} and several consequences from Wilson~\cite{Wilson-1967}.

\begin{theorem}\label{T-1-7}
Let $M$ be a smooth manifold of dimension $d\geq 1$ that admits a Lyapunov-Reeb function $f$. The following properties hold:
\begin{itemize}
\item[$(i)$] Given that $M$ is equipped with a Riemannian metric, the negative gradient flow of $f$ starting at any point in $M$ exists for all $t\geq 0$ and converges to the critical point of $f$ as $t\to+\infty$.
\item[$(ii)$] $M$ is diffeomorphic to $\mathbb R^d$.
\item[$(iii)$] For any $c>\inf f$, the set $M_c=\big\{x\in M: f(x)<c\big\}$ is diffeomorphic to $\mathbb R^d$, and the set $L_c=\big\{x\in M: f(x)=c\big\}$ is homeomorphic to $S^{d-1}$.
\end{itemize}
\end{theorem}

\begin{remark}
The definition of L-R functions is motivated by the two relations: $(1)$ Up to a constant, an L-R function serves as a Lyapunov function for its negative gradient flow (see Remark~\ref{R-2-4}); $(2)$ As in the Reeb Sphere Theorem~\cite{Milnor-1964}, the critical point of an L-R function is allowed to be degenerate.
\end{remark}

To prove Theorems~\ref{T-1-3} and~\ref{T-1-4}, it suffices to construct L-R functions on the relevant spaces. We first introduce some concepts motivated by the work of Cantarella, Demaine, Iben and O'Brien~\cite{Cantarella-2004}.

\begin{definition}\label{D-1-9}
For $U=\mathcal M(R_m,\ell)$ or $U=\mathcal M^{+}(C_m,\ell)$,  a proper smooth map $\Phi:U\to[0,+\infty)$ is called a \textbf{strain energy} on $U$ if it satisfies
\[
\frac{\mathrm{d}\Phi(p_t)}{\mathrm{d}t}<0
\]
for every expansive motion $\{p_t\}_{-\varepsilon<t<\varepsilon}$ ($\varepsilon>0$) in $U$ and every $t\in(-\varepsilon,\varepsilon)$.
\end{definition}

Here a smooth family $\{p_t\}_{-\varepsilon<t<\varepsilon}$ ($\varepsilon>0$) of linkages in $U$ is called an \textbf{expansive motion} in $U$  if for all $t\in (-\varepsilon,\varepsilon)$ and
all pairs of vertices $v_i,v_j\in V$,
\[
\frac{\mathrm{d}}{\mathrm{d}t}
~\| p_{t}(v_i)-p_{t}(v_j)\|\geq 0,
\]
with at least one of the inequalities being strict.

\begin{remark}\label{R-1-10}
Let $m\geq 3$. It is easy to see that the function
\[
\Phi(p)=\sum_{\substack{[v_i,v_j]\in E \\ v_k\in V\setminus\{v_i,v_j\}}}\frac{1}
{\left(\|p(v_i)-p(v_k)\|+\|p(v_j)-p(v_k)\|-\|p(v_i)-p(v_j)\|\right)^2}
\]
introduced by Cantarella et al.~\cite{Cantarella-2004} is a strain energy on
$\mathcal M(R_m,\ell)$ (resp. $\mathcal M^{+}(C_m,\ell)$).
\end{remark}

We find that strain energies on $\mathcal M(R_m,\ell)$ serve as L-R functions.

\begin{theorem}\label{T-1-11}
Let $m\geq 2$ and let $\ell:E\to\mathbb R_{+}$ be an arbitrary assignment. Then every strain energy on $\mathcal M(R_m,\ell)$ is a Lyapunov-Reeb function on
$\mathcal M(R_m,\ell)$.
\end{theorem}

Next we construct L-R functions on $\mathcal M^{+}(C_m,\ell)$. To this end, we introduce two functions on $\mathcal M^{+}(C_m)$ (and hence on $\mathcal M^{+}(C_m,\ell)$).
Recall that each $p\in \mathcal M^{+}(C_m)$ gives a positively oriented simple polygon $p(C_m)$. Let $A(p)$ be the area enclosed by $p(C_m)$. This gives rise to a function
\[
A: \mathcal M^{+}(C_m)\to \mathbb R_{+}.
\]
Following Shimamoto and Wootters~\cite{Shimamoto-2014}, for $i=1,2,\cdots,m$, we set
%\[
%w_i(p)= %\begin{cases}~\exp\left(\frac{1}{\pi-\alpha_i(p)}\right), & %\, \alpha_i(p)>\pi, \\ ~0, &\, \alpha_i(p) \leq %\pi,\end{cases}
%\]
\[
w_i(p)=
\begin{cases}
\begin{aligned}
&\exp\left(\frac{1}{\pi-\alpha_i(p)}\right), &&\alpha_i(p)>\pi,\\
&0, &&\alpha_i(p)\leq\pi,
\end{aligned}
\end{cases}
\]
where $\alpha_i(p)$ denotes the interior angle of $p(C_m)$ at vertex $p(v_i)$. We then define a function $w: \mathcal M^{+}(C_m)\to [0,1)$ as
\[
w=\frac{1}{m}\left(\sum\nolimits_{i=1}^m w_i\right).
\]
Note that $w(p)=0$ if and only if $p\in \mathcal M^{+}(C_m)$ corresponds to a convex polygon. We acquire L-R functions on $\mathcal M^{+}(C_m,\ell)$ as follows.

\begin{theorem}\label{T-1-12}
Let $m\geq 3$ and let $\ell:E\to\mathbb R_{+}$ satisfy condition~$\mathrm{\mathbf{(c1)}}$. Suppose $\Phi:\mathcal M^{+}(C_m,\ell)\to [0,+\infty)$ is a strain energy. Then
\[
f=\frac{1}{A}+w\Phi
\]
is a Lyapunov-Reeb function on $\mathcal M^{+}(C_m,\ell)$.
\end{theorem}

The proofs of the above two theorems rely on a fundamental result established by Connelly, Demaine and Rote~\cite{Connelly-2003} (see also the work of Rote, Santos and Streinu~\cite{Rote-2003}), which states that every non-straight arm linkage and every non-convex cycle linkage admits expansive motions. Specifically, this property implies that strain energies have no critical points in certain subsets of moduli spaces.

By incorporating a projection method, we further construct L-R functions on $\mathcal M(R_m)$ and $\mathcal M^{+}(C_m)$ (see Theorems~\ref{T-3-5} and~\ref{T-4-16}). Combining these derivations with Theorem~\ref{T-1-7}, we then establish Theorems~\ref{T-1-3} and~\ref{T-1-4}. Notably, this framework yields a number of additional corollaries. Here we present two results which describe foliation structures of the moduli spaces via level sets of L-R functions.

\begin{corollary}\label{C-1-13}
Let $f$ be a Lyapunov-Reeb function on $\mathcal M(R_m,\ell)$ with $m\geq 3$. Given $c>\inf f$, $f^{-1}(c)$ is homeomorphic to $S^{m-3}$.
\end{corollary}

\begin{corollary}\label{C-1-14}
Suppose $\ell:E\to\mathbb R_{+}$ satisfies condition~$\mathrm{\mathbf{(c1)}}$ and
suppose $f$ is a Lyapunov-Reeb function on $\mathcal M^{+}(C_m,\ell)$ with $m\geq4$. Given $c>\inf f$, $f^{-1}(c)$ is homeomorphic to $S^{m-4}$.
\end{corollary}


\begin{remark}
Corollaries~\ref{C-1-13} and~\ref{C-1-14} settle the following strengthened version of Carpenter's Rule Problem: \emph{Under what conditions can a polygonal linkage be continuously deformed to any other linkage in the same level set of an L-R function, such that during the deformation process no self-intersection is created and the L-R function remains constant?}
\end{remark}

Other applications of L-R functions include devising algorithms for the original Carpenter's Rule Problem through associated negative gradient flows. In fact, L-R functions shed light on a broader question: \emph{How can one generate an animation from two images?} More precisely, we will introduce a renormalization algorithm for the Linkage (Configuration) Refolding Problem~\cite{Cantarella-2004,Demaine-2009}: \emph{Given two polygonal linkages (configurations), construct an ``optimal'' motion connecting them.}
Further details can be found in~\S\ref{S-5}.

The paper is organized as follows: In~\S\ref{S-2} we develop Morse theory for open manifolds and prove Theorem~\ref{T-1-7}. In~\S\ref{S-3} we construct L-R functions on the moduli spaces of arm linkages and configurations, and establish Theorems~\ref{T-1-11},~\ref{T-1-3} and Corollary~\ref{C-1-13}. In~\S\ref{S-4} we study the moduli spaces of cycle linkages and configurations, deducing
Theorems~\ref{T-1-12},~\ref{T-1-4} and Corollary~\ref{C-1-14}. In~\S\ref{S-5} we present algorithms addressing the Carpenter's Rule Problem and the Linkage (Configuration) Refolding Problem. In~\S\ref{S-6} we propose several problems for further investigation.

Throughout this paper, a manifold is assumed to be a second-countable Hausdorff space that is locally homeomorphic to Euclidean space, following the standard terminology.

\section{Morse theory in the setting of open manifolds}\label{S-2}

This section is dedicated to showing Theorem~\ref{T-1-7}. Guided by the spirit of  the proof of the Reeb Sphere Theorem~\cite{Reeb-1946,Milnor-1964}, we employ negative gradient flows along with certain modifications. In particular, we draw on Wilson's insight~\cite{Wilson-1967} that a modified negative gradient flow can induce a homotopy equivalence from the given level set to a specific sphere. Additionally, we shall make use of a powerful tool\textemdash the Generalized Poincar\'{e} Conjecture (see, e.g.,~\cite{Smale-1961,Freedman-1982,Perelman-2002,Perelman-2003.3,Perelman-2003.7}).


\subsection{Negative gradient flows}

First we introduce some concepts and results on negative gradient flows. For an exposition, see the textbook by Lee~\cite{Lee-2018}.

Let $M$ be a smooth manifold. A flow on $M$ is a smooth map $\varphi:\mathcal D\to M$ with the following properties:
\begin{itemize}
\item[$(1)$]$\mathcal D$ is an open subset of $\mathbb R\times M$ such that $\{0\}\times M\subset \mathcal D$;
\item[$(2)$] $\varphi(0,p)=p$;
\item[$(3)$] If $(s,p)\in\mathcal D$ and $\big(t,\varphi(s,p)\big)\in\mathcal D$ such that $(t+s,p)\in\mathcal D$, then
\[
\varphi\big(t,\varphi(s,p)\big)=\varphi(t+s,p).
\]
\end{itemize}
A flow $\varphi:\mathcal D\to M$ is said to be \textbf{forward-complete} (resp. \textbf{backward-complete}) if
\[
[0,+\infty)\times M\subset \mathcal D\quad \left(\text{resp}.\;\;(-\infty,0]\times M\subset \mathcal D\right),
\]
and it is \textbf{complete} if $\mathcal D=\mathbb R\times M$. For a complete flow $\varphi:\mathbb R\times M\to M$, note that the flow map $\varphi(t,\cdot):M\to M$ is a diffeomorphism for every $t\in \mathbb R$.

A flow $\varphi:\mathcal D\to M$ is said to be generated by a smooth vector field $X$ on $M$ if it satisfies
\[
\frac{\mathrm{d}\varphi(t,p)}{\mathrm{d}t}
=X\big(\varphi(t,p)\big).
\]

We recall the Fundamental Theorem on Flows (see, e.g.,~\cite[Theorem 9.12]{Lee-2018}).

\begin{theorem}\label{T-2-1}
Every smooth vector field $X$ on $M$ generates a unique maximal flow on $M$. Moreover, the resulting flow is complete if $X$ is compactly supported.
\end{theorem}

We mention that a \textbf{maximal flow} means it admits no extension to a flow on a larger domain, and a vector field $X$ is \textbf{compactly supported} if its support,  given by
$\operatorname{supp} X:=\overline{\left\{p\in M: X(p)\neq 0\right\}}$, is compact.

Suppose $M$ is endowed with a Riemannian metric. The gradient of a smooth function $f$ on $M$ is the vector field $\nabla f$ such that
\[
\langle \nabla f, X\rangle= X(f)
\]
holds for all smooth vector fields $X$ on $M$. To develop Morse theory,  we need to study the negative gradient flow of $f$, i.e., the maximal flow that satisfies
\[
\frac{\mathrm{d}\varphi(t,p)}{\mathrm{d}t}
=-\nabla f\big(\varphi(t,p)\big).
\]

For open manifolds, the negative gradient may not be compactly supported, and the induced flow may fail to be complete.

\begin{example}
Define a function on $\mathbb R^2$ by $f(x,y)=x^2+y^4$. The negative gradient flow of $f$ is expressed as
\[
\varphi\big(t,(x,y)\big)=\left(e^{-2t}x,\frac{y}{(1+8ty^2)^{\frac{1}{2}}}\right).
\]
This flow is not complete. More specifically, the flow map $\varphi(t,\cdot):\mathbb R^2\to \mathbb R^2$ is non-surjective for $t>0$ and not well defined for $t<0$.
\end{example}


Fortunately, forward-completeness holds for the negative gradient flows of all proper, smooth, and lower-bounded functions.

\begin{lemma}\label{L-2-2}
Let $M$ be a smooth Riemannian manifold, and let $f:M\to\mathbb R$ be a proper, smooth, and lower-bounded function. Then the negative gradient flow of $f$ is forward-complete.
\end{lemma}
\begin{proof}
In light of Theorem~\ref{T-2-1}, there exists a unique maximal flow $\varphi:\mathcal D\to M$ generated by $-\nabla f$. For each $p\in M$, let $J_p$ be the maximal interval such that
\[
J_p\times\{p\}\subset \mathcal D.
\]
It suffices to show $J_p$ has no upper bounds. Assume the contrary, then $\beta:=\sup J_p$ is finite. Define the integral curve starting at $p$ as
\[
\begin{aligned}
\gamma_p\;:\;[0,\beta)\;&\longrightarrow\; M\\
\;\;\;t\quad\;&\longmapsto\;\varphi(t,p).
\end{aligned}
\]
By the Escape Lemma~\cite[Lemma 9.19]{Lee-2018},  $\gamma_p\big([0,\beta)\big)$ cannot be contained in any compact subset of $M$.

On the other hand, a direct calculation gives
\[
\frac{\mathrm{d}f\big(\gamma_p(t)\big)}{\mathrm{d}t}
=-\| \nabla f \|^2\leq 0,
\]
which yields
\[
f\big(\gamma_p(t)\big)\leq f\big(\gamma_p(0)\big)=f(p).
\]
Recall that $f$ is lower-bounded. We obtain
\[
\gamma_p\big([0,\beta)\big)\subset f^{-1}\big([a,b]\big),
\]
where $a=\min f$ and $b=f(p)$. By the properness of $f$, the preimage $f^{-1}\big([a,b]\big)$ is compact, which contradicts the fact that $\gamma_p\big([0,\beta)\big)$ cannot be contained in any compact subset of $M$. We thus derive $\beta=+\infty$ and conclude the lemma.
\end{proof}

In addition, one asserts the convergence of the flow assuming further that $f$ possesses exactly one critical point.
\begin{theorem}\label{T-2-3}
Let $M$ be a smooth Riemannian manifold, and let
$f: M\to\mathbb R$ be a Lyapunov-Reeb function with $q$ being the critical point. Suppose $\varphi:\mathcal D\to M$ is the negative gradient flow of $f$. Then for any $p\in M$,
\[
\lim_{t\to+\infty}\varphi(t,p)=q.
\]
\end{theorem}

\begin{proof}
Consider the integral curve
\[
\begin{aligned}
\gamma_p\;:\;[0,+\infty)\;&\longrightarrow\; M\\
\;\;\;t\quad\;&\longmapsto\;\varphi(t,p).
\end{aligned}
\]
As in Lemma~\ref{L-2-2}, it is easy to see  $f\big(\gamma_p(t)\big)$ is a decreasing function of $t$. Since $f$ is lower-bounded,   $\lim_{t\to+\infty}f\big(\gamma_p(t)\big)$ exists. For each positive integer $n$, Lagrange's Mean Value Theorem implies there exists $t_n\in(n,n+1)$ such that
\[
f\big(\gamma_p(n+1)\big)-f\big(\gamma_p(n)\big)
=\frac{\mathrm{d}f\big(\gamma_p(t)\big)}{\mathrm{d}t}\bigg|_{t=t_n}
=-\|\, \nabla f\big(\gamma_p(t_n)\big)\,\|^2.
\]
Because $\lim_{t\to+\infty}f\big(\gamma_p(t)\big)$ exists, we have
\[
f\big(\gamma_p(n+1)\big)-f\big(\gamma_p(n)\big)\to 0
\]
as $n\to+\infty$. Consequently,
\[
\|\,\nabla f\big(\gamma_p(t_n)\big)\,\|^2\to 0.
\]
Note that $\big\{\gamma_p(t_n)\big\}$ lies in the compact set $f^{-1}\big([a,b]\big)$,
where
\[
a=\min f=f(q)\quad\text{and}\quad b=f(p).
\]
By passing to a subsequence if necessary, we may assume $\big\{\gamma_p(t_n)\big\}$ converges to some $q^\ast\in M$. Then $\nabla f (q^\ast)=0$, which implies $q^\ast$ is a critical point of $f$. Given that $f$ has only one critical point, we deduce $q^\ast=q$, meaning
\[
\gamma_p(t_n)\to q.
\]
As a result,
\[
\lim_{t\to+\infty}f\big(\gamma_p(t)\big)
=\lim_{n\to+\infty}f\big(\gamma_p(t_n)\big)=f(q)=\min f.
\]
We claim that $\gamma_{p}(t)$ converges to $q$. Otherwise, there exists a subsequence $\big\{\gamma_{p}(s_k)\big\}$ converging to some $\tilde{q}\neq q$. It follows that
\[
\lim_{k\to+\infty}f\big(\gamma_p(s_k)\big)=f\big(\tilde{q}\big)>\min f,
\]
contradicting our earlier conclusion. Hence
$\varphi(t,p)=\gamma_p(t)\to q$ as $t\to+\infty$.
\end{proof}

\begin{remark}\label{R-2-4}
Set $h=f-f(q)$. Then $h$ serves as a global Lyapunov function for its negative gradient flow. By virtue of the Barba\v{s}in-Krasovski\v{i} Theorem~\cite{Barba-1954},  $q$ is a globally asymptotically stable equilibrium point of the system, implying each point is attracted to $q$ by the flow. This offers an alternative approach to Theorem~\ref{T-2-3}.
\end{remark}

\subsection{Morse theory on open manifolds}

To streamline the arguments, we now restate property $(ii)$ of Theorem~\ref{T-1-7} as the following result. As mentioned before, this can be viewed as an analog of the Reeb Sphere Theorem~\cite{Reeb-1946,Milnor-1964} in the setting of open manifolds.

\begin{theorem}\label{T-2-5}
Suppose $M$ is a smooth manifold of dimension $d\geq1$ that admits a Lyapunov-Reeb function $f$. Then $M$ is diffeomorphic to $\mathbb R^{d}$.
\end{theorem}

Before giving the proof, we introduce two preliminary lemmas. The first one is due to Brown~\cite{Brown-1961} and Stallings~\cite{Stallings-1962}. The following version together with a proof can be seen in the lecture note of Milnor~\cite[p.~168]{Milnor-1964}.
\begin{lemma}[Brown-Stallings]\label{L-2-6}
Let $M$ be a smooth manifold such that every compact subset is contained in an open subset that is diffeomorphic to $\mathbb{R}^d$. Then $M$ is diffeomorphic to $\mathbb{R}^d$.
\end{lemma}

\begin{remark}
Milnor's original statement~\cite[p.~168]{Milnor-1964} assumes $M$ is paracompact. This can be derived from our presupposition that manifolds are second-countable Hausdorff spaces. For a proof,  see, e.g.,~\cite[Theorem 1.15]{Lee-2018}.
\end{remark}


The next lemma addresses the existence of smooth bump functions.
\begin{lemma}\label{L-2-8}
Given $a,b\in \mathbb R$ with $a<b$, there exists a smooth function $\eta_{a,b}:\mathbb R\to\mathbb R$ such that $\eta_{a,b}\equiv1$ on $(-\infty,a]$, $\eta_{a,b}\equiv 0$ on $[b,+\infty)$, and $0<\eta_{a,b}<1$ on $(a,b)$.
\end{lemma}
\begin{proof}
It is easy to see that the function
\[
\eta_{a,b}(x)=
\begin{cases}
\begin{aligned}
&1, &&x\leq a,\\
&\exp\left(\dfrac{x-a}{x-b}\right), &&a<x<b,\\
&0, &&x\geq b
\end{aligned}
\end{cases}
\]
satisfies the required properties.
\end{proof}

We are ready to demonstrate Theorem~\ref{T-2-5}.

\begin{proof}[\textbf{Proof of Theorem~\ref{T-2-5}}]
By Lemma~\ref{L-2-6}, it suffices to show that every compact subset $K\subset M$ is contained in an open subset $W\subset M$ that is diffeomorphic to $\mathbb{R}^d$.

For this purpose, we first equip $M$ with a Riemannian metric. Let $q$ be the unique critical point of $f$, and let $B_q\subset M$ be a relatively compact neighborhood of $q$ that is diffeomorphic to $\mathbb{R}^d$. Suppose $\varphi:\mathcal D\to M$ is the negative gradient flow of $f$. By virtue of Proposition~\ref{P-2-10} below, there exists $T>0$ such that
\[
\varphi(T,K)\subset B_q.
\]
Let $\eta_{a,b}: \mathbb R\to\mathbb R$ be given as in Lemma~\ref{L-2-8} with
\[
a=1+\max\nolimits_{p\in K} f(p)\quad\text{and}\quad b=2+\max\nolimits_{p\in K} f(p).
\]
Then $-\eta_{a,b}(f)\nabla f$ is a smooth vector field on $M$ whose support equals
\[
K_b:=\big\{p\in M: f(p)\leq b\big\}.
\]
Since $f: M\to\mathbb R$ is proper and lower-bounded,  $K_b$ is clearly a compact set. By Theorem~\ref{T-2-1}, $-\eta_{a,b}(f)\nabla f$ generates a complete flow
$\varphi_\eta:\mathbb R\times M\to M$. Thus for every
$t\in \mathbb R$, the flow map
$\varphi_\eta(t,\cdot):M \to M$ is a diffeomorphism. Note that
\[
\frac{\mathrm{d}f\big(\varphi_\eta(t,p)\big)}{\mathrm{d}t}
=-\eta_{a,b}(f)\,\| \nabla f \|^2\leq 0,
\]
and $\eta_{a,b}(f)\equiv1$ on
$K_a:=\big\{p\in M: f(p)\leq a\big\}$.
For $t\geq 0$ and $p\in K_a$, we derive
 \[
\varphi_\eta(t,p)=\varphi(t,p)\in K_a.
\]
Because $T>0$ and $K\subset K_a$, this implies
\[
\varphi_\eta(T,K)=\varphi(T,K)\subset B_q.
\]
Setting $W=\varphi_\eta(-T,B_q)$ gives the required open set, proving the assertion.
\end{proof}

\begin{remark}
As previously noted, the original flow $\varphi$ may not be complete, which impedes defining the backward flow map on $B_q$. Fortunately, this issue is remedied by the modified flow $\varphi_\eta$, allowing the required construction to proceed.
\end{remark}


The following proposition provides a uniform bound for the ``arrival time'' of the negative gradient flow when the initial points lie in a compact set.
\begin{proposition}\label{P-2-10}
There exists $T>0$ such that for all $t\geq T$,
\[
\varphi(t,K)\subset B_q.
\]
\end{proposition}

\begin{proof}
Define
\[
\mu=\min_{p\in\partial B_q} f(p)\quad\text{and}\quad M_{\mu}=\big\{p\in M: f(p)<\mu\big\}.
\]
We claim $M_{\mu}\subset B_q$. Otherwise, there exists another critical point $\tilde{q}$ of $f$ such that
\[
f(\tilde{q})=\min_{p\in M\setminus B_q} f(p)<\mu=\min_{p\in\partial B_q} f(p),
\]
contradicting the assumption that $f$ has exactly one critical point.
For every $p\in M$, Theorem~\ref{T-2-3} guarantees the existence of $t_p\geq 0$ with the property that
\[
\varphi(t_p,p)\in M_{\mu}.
\]
Since $\varphi(t_p,\cdot): M\to M$ is a smooth map and $M_{\mu}$ is an open subset, one can find a neighborhood $U_p\subset M$ of $p$ satisfying
\[
\varphi(t_p,U_p)\subset M_{\mu}.
\]
Note that $\big\{U_p\big\}_{p\in K}$ forms an open cover of the compact set $K$ and $f\big(\varphi(t,p)\big)$ is a decreasing function of $t$. By selecting a finite subcover, we readily obtain $T>0$ such that for all $t\geq T$,
\[
\varphi(t,K)\subset M_{\mu}\subset B_q,
\]
thereby establishing the proposition.
\end{proof}

To derive Theorem~\ref{T-1-7}, we invoke the following result affirming the Generalized Poincar\'{e} Conjecture for differentiable manifolds. It was proved by Smale~\cite{Smale-1961} in dimensions greater than $4$, by Freedman~\cite{Freedman-1982} in dimension $4$, and by Perelman~\cite{Perelman-2002,Perelman-2003.3,Perelman-2003.7} in dimension $3$. For dimensions lower than 3, it has long been known by the classification of manifolds in those dimensions.

\begin{theorem}[Smale, Freedman, Perelman]\label{T-2-11}
Every closed smooth $d$-manifold that has the homotopy type of $S^d$ is homeomorphic to $S^d$.
\end{theorem}

\begin{remark}
The smoothness condition in the above theorem can actually be dropped. For $d\geq5$, this was established by Newman~\cite{Newman-1966} building on research of Stallings~\cite{Stallings-1960} and Zeeman~\cite{Zeeman-1961}. For $d=4$ and $d=3$, the result is attributed to Freedman~\cite{Freedman-1982} and Perelman~\cite{Perelman-2002,Perelman-2003.3,Perelman-2003.7}, respectively.
\end{remark}


We are now in a position to show Theorem~\ref{T-1-7}.
\begin{proof}[\textbf{Proof of Theorem~\ref{T-1-7}}]
Property $(i)$ follows straightforwardly from Lemma~\ref{L-2-2} and Theorem~\ref{T-2-3}, while property $(ii)$ was established in Theorem~\ref{T-2-5}.

To address property $(iii)$, we first note that $M_c=\big\{x\in M: f(x)<c\big\}$ is a smooth manifold of dimension $d$. Define
\[
h=\frac{1}{c-f},
\]
which serves as a Lyapunov-Reeb function on $M_c$. By property $(ii)$ (equivalently, Theorem~\ref{T-2-5}), we infer that $M_c$ is diffeomorphic to $\mathbb{R}^d$.

It remains to characterize the topology of the level set
$L_c=\big\{x\in M: f(x)=c\big\}$.  Since $f$ is proper and $c>\inf f$ is a regular value, the Regular Value Theorem~\cite[Corollary 5.14]{Lee-2018} implies $L_c$ is a closed $(d-1)$-dimensional smooth manifold. In view of Theorem~\ref{T-2-11}, to prove $L_c$
is homeomorphic to $S^{d-1}$, we need to show $L_c$ has the homotopy type of $S^{d-1}$.

To this end, we refine a technique by Wilson~\cite{Wilson-1967} to construct a homotopy equivalence between $L_c$ and $M\setminus\{q\}$, where $q$ is the unique critical point of $f$. To start, for any $p\in M\setminus\{q\}$, consider the maximal integral curve $\xi_p:I_p\to M\setminus\{q\}$
satisfying the ordinary differential equation
\[
\frac{\mathrm d \xi_p}{ \mathrm d t} = -\frac{\nabla f}{\|\nabla f\|^2}
\]
with initial condition $\xi_p(0)=p$. A routine calculation yields
\[
\frac{\mathrm d f\big(\xi_{p}(t)\big)}{\mathrm d t}\equiv-1.
\]
By the Escape Lemma~\cite[Lemma 9.19]{Lee-2018}, the maximal interval is given by
\[
I_p=\big(-\infty,f(p)-f(q)\big).
\]
Letting $p$ vary over $L_c$ defines a smooth map
\[
\begin{aligned}
\psi:\; \mathcal D_c\;\,&\longrightarrow\; M\setminus\{q\}\\
          (t,p)\;&\longmapsto\;\; \xi_p(t),
\end{aligned}
\]
where $\mathcal D_c=\big(-\infty,c-f(q)\big)\times L_c$. By Proposition~\ref{P-2-13} below,  $\psi$ is a homeomorphism. Since $M$ is diffeomorphic to $\mathbb{R}^d$,  both $M\setminus\{q\}$ and $\mathcal D_c$ have the homotopy type of $S^{d-1}$. Noting that $L_c$ is homotopy equivalent to $\mathcal D_c$, we prove $L_c$ is a closed smooth manifold homotopy equivalent to $S^{d-1}$. Overall, these topological relations are summarized as follows:
\[
L_c\simeq \mathcal D_c\cong M\setminus\{q\}\cong\mathbb R^d\setminus\{\mathbf{0}\}\simeq S^{d-1},
\]
where $\mathbf{0}\in \mathbb R^d$ denotes the origin, and the symbols ``$\simeq$'' and ``$\cong$'' denote homotopy equivalence and homeomorphism between topological spaces, respectively. In light of Theorem~\ref{T-2-11}, we conclude that $L_c$ is homeomorphic to $S^{d-1}$, as desired.
\end{proof}





\begin{proposition}\label{P-2-13}
The map $\psi:\mathcal D_c\to M\setminus\{q\}$ is a homeomorphism.
\end{proposition}
\begin{proof}
First it is straightforward to see that
\[
\dim\left(\mathcal D_c\right)=\dim \left(M\setminus\{q\}\right)=d.
\]
For $d=1$, $\psi$ is clearly a homeomorphism. Henceforth, we assume $d\geq 2$ and observe the following properties:
\begin{itemize}
\item[\textbf{(x1)}] $\psi$ is proper. This follows from the Escape Lemma~\cite[Lemma 9.19]{Lee-2018}.
\item[\textbf{(x2)}] $\psi$ is injective. Suppose $(t_1,p_1),(t_2,p_2)\in\mathcal D_c$ and $y\in M\setminus\{q\}$ satisfy
    \[
\xi_{p_1}(t_1)=\xi_{p_2}(t_2)=y.
    \]
    We need to show $(t_1,p_1)=(t_2,p_2)$. As before, a direct calculation gives
    \[
    \frac{\mathrm d f\big(\xi_{p_1}(t)\big)}{\mathrm d t}\equiv-1,
    \]
    which yields
    \[
    -t_1=f\big(\xi_{p_1}(t_1)\big)-f\big(\xi_{p_1}(0)\big)=f(y)-f(p_1).
    \]
    Similarly,
    \[
    -t_2=f(y)-f(p_2).
    \]
    Because $p_1,p_2\in L_c$, we have
    \[
    t_1=f(p_1)-f(y)=c-f(y)=f(p_2)-f(y)=t_2.
    \]
Moreover,
    \[
    p_1=\xi_y(-t_1)=\xi_y(-t_2)=p_2.
    \]
Therefore, $(t_1,p_1)=(t_2,p_2)$, confirming that $\psi$ is injective.
\end{itemize}

From property~\textbf{(x1)}, $\psi(\mathcal D_c)$ is a closed subset of $M\setminus\{q\}$. Meanwhile, by Brouwer's Invariance of Domain Theorem,
property~\textbf{(x2)} implies $\psi(\mathcal D_c)$ is open in $M\setminus\{q\}$. Since $\psi(\mathcal D_c)$ is clopen in $M\setminus\{q\}$ and  $M\setminus\{q\}$ is connected for $d\geq 2$, we deduce that
$\psi(\mathcal D_c)=M\setminus\{q\}$. In summary, $\psi$ is a bijective smooth map and thus a homeomorphism.
\end{proof}

\begin{remark}
Wilson~\cite{Wilson-1967} obtained similar results to properties $(ii)$ and $(iii)$ of Theorem~\ref{T-1-7} under stronger assumptions that $M$ is an open subset of
$\mathbb R^d$ and $f$ is a global Lyapunov function for a complete flow on $M$. By comparison, after relaxing Wilson's conditions (particularly the flow's completeness), our theorem offers broader applicability  in characterizing the topology of smooth manifolds.
\end{remark}



\section{Moduli spaces of arm linkages and configurations}\label{S-3}
In this section, we determine the topology of the moduli spaces of arm linkages and configurations. As noted earlier, our strategy is to construct L-R functions on these spaces. To achieve this goal, we need to examine the smooth structures of $\mathcal M(R_m)$ and $\mathcal M(R_m,\ell)$.

Let us begin with the moduli space $\mathcal M(R_m)$. For each $p\in \mathcal M(R_m)$, via an appropriate orientation-preserving isometry, we may assume
\[
\begin{aligned}
&p(v_1)=(0,0),\quad p(v_2)=(\rho_1,0),\\
&p(v_3)=p(v_2)+(\rho_2\cos\theta_2,\rho_2\sin\theta_2),\\
&\quad\qquad\qquad\qquad\vdots\\
&p(v_{m})=p(v_{m-1})+(\rho_{m-1}\cos\theta_{m-1},\rho_{m-1}\sin\theta_{m-1}),
\end{aligned}
\]
where $(\rho,\theta)
:=(\rho_1,\cdots,\rho_{m-1},\theta_2,\cdots,\theta_{m-1})\in \mathbb R_{+}^{m-1}\times\left(\mathbb R/2\pi\mathbb Z\right)^{m-2}$.
Noting that each
$p\in \mathcal M(R_m)$ is uniquely determined by $(\rho,\theta)$, we can identify $\mathcal M(R_m)$ with a subset of $\mathbb R_{+}^{m-1}\times\left(\mathbb R/2\pi\mathbb Z\right)^{m-2}$. Since the self-avoiding property is stable under small perturbations,  $\mathcal M(R_m)$ is actually an open subset of
$\mathbb R_{+}^{m-1}\times\left(\mathbb R/2\pi\mathbb Z\right)^{m-2}$, thus inheriting the induced smooth structure from the ambient space. Meanwhile,
\[
\mathcal M(R_m,\ell)=\left\{(\rho,\theta)\in\mathcal M(R_m):\rho=\vec{\ell}\right\},
\]
where $\vec{\ell}=\big(\ell([v_1,v_2]),\cdots,\ell([v_{m-1},v_m])\big)$ denotes the length vector. This implies that $\mathcal M(R_m,\ell)$ carries the smooth structure of an embedded submanifold of $\mathcal M(R_m)$.

In summary, we obtain the following lemma, part of which has appeared in the works of Shimamoto and Vanderwaart~\cite{Shimamoto-2005}
 and Farber~\cite[Chap. 3]{Farber-2008}.
\begin{lemma}\label{L-3-1}
Let $m\geq 2$ and let $\ell:E\to\mathbb R_{+}$ be an arbitrary assignment. The following statements hold:
\begin{itemize}
\item[$(i)$]$\mathcal M(R_m)$ is a smooth manifold of dimension $2m-3$.
\item[$(ii)$]$\mathcal M(R_m,\ell)\subset \mathcal M(R_m)$ is a smooth submanifold of dimension $m-2$.
\end{itemize}
\end{lemma}

To establish that strain energies on $\mathcal M(R_m,\ell)$ are L-R functions, we invoke the following result by Connelly, Demaine and Rote~\cite[Theorem 3]{Connelly-2003}, which played a pivotal role in resolving the Carpenter's Rule Problem. For alternative proofs, see also the works of Rote, Santos and Streinu~\cite[Theorem 4.3]{Rote-2003} and Farber~\cite[Chap.~3]{Farber-2008}. We call an element in $\mathcal M(R_m,\ell)$ (resp. $\mathcal M(R_m)$) a \textbf{straight linkage} (resp.~\textbf{straight configuration}) if all its vertices are collinear.

\begin{theorem}[Connelly-Demaine-Rote]\label{T-3-2}
Suppose $p\in \mathcal M(R_m,\ell)$ is not the straight linkage. Then there exists an expansive motion $\{p_t\}_{-\varepsilon<t<\varepsilon}$ ($\varepsilon>0$) in $\mathcal M(R_m,\ell)$ such that $p_0=p$.
\end{theorem}


We proceed to show Theorem~\ref{T-1-11}.

\begin{proof}[\textbf{Proof of Theorem \ref{T-1-11}}]
Suppose $\Phi:\mathcal M(R_m,\ell)\to[0,+\infty)$ is a strain energy. It is easy to see that $\Phi$ has a minimum point $q\in \mathcal M(R_m,\ell)$. We shall prove that $q$ is the unique critical point of $\Phi$, which confirms $\Phi$ is a Lyapunov-Reeb function.

First we claim $q$ is the straight linkage. If this is not the case, then Theorem~\ref{T-3-2} implies there exists an expansive motion $\{p_t\}_{-\varepsilon<t<\varepsilon}$ ($\varepsilon>0$) in
$\mathcal M(R_m,\ell)$ such that $p_0=q$. Since $\Phi$ is a strain energy, we have
\[
\frac{\mathrm{d}\Phi(p_t)}{\mathrm{d}t}\bigg|_{t=0}<0,
\]
contradicting the fact that $p_0=q$ is a minimum point of $\Phi$.

Next we verify that $\Phi$ has no critical points other than $q$. Assume, for contradiction, that $q'$ is another critical point of $\Phi$. Then $q'$ cannot be the straight linkage. Applying Theorem~\ref{T-3-2} again, we obtain an expansive motion $\{p'_t\}_{-\varepsilon'<t<\varepsilon'}$ ($\varepsilon'>0$) in $\mathcal M(R_m,\ell)$ with $p'_0=q'$. Because $q'$ is a critical point of $\Phi$, a straightforward computation yields
\[
\frac{\mathrm{d}\Phi(p'_t)}{\mathrm{d}t}\bigg|_{t=0}=0,
\]
which contradicts the definition of strain energies. Thus the theorem is proved.
\end{proof}


To construct L-R functions on $\mathcal M(R_m)$, we devise a projection method. Precisely, for each $p\in \mathcal M(R_m)$, let $\varsigma(p)\in \mathcal M(R_m)$ be the straight configuration sharing the same length vector as $p$. This induces the following projection map
\[
\begin{aligned}
\varsigma\;:\;\mathcal M(R_m)\;&\longrightarrow\;\mathcal L(R_m)\\
p\;\;\;\,&\longmapsto\;\;\varsigma(p),
\end{aligned}
\]
where $\mathcal L(R_m)\subset\mathcal M(R_m)$ denotes the subspace of straight configurations. In terms of coordinates, $\mathcal L(R_m)$ is characterized by
\[
\mathcal L(R_m)=\left\{(\rho,\theta)\in\mathcal M(R_m):\theta=\vec{0}\right\}.
\]
Here $\vec{0}:=(0,\cdots,0)\in \left(\mathbb R/2\pi\mathbb Z\right)^{m-2}$ represents the null vector. Obviously, $\mathcal L(R_m)$ is a smooth manifold diffeomorphic to $\mathbb R_{+}^{m-1}$.

\begin{definition}
A smooth map $\Phi:\mathcal M(R_m)\to [0,+\infty)$ is called a
\textbf{$\varsigma$-potential} on $\mathcal M(R_m)$ if the following properties hold:
\begin{itemize}
\item[\textbf{(a1)}] For every $q\in \mathcal L(R_m)$, the restriction of $\Phi$ to $\varsigma^{-1}(q)$ is a strain energy;
\item[\textbf{(a2)}]For every compact set $K\subset \mathcal L(R_m)$, the restriction of $\Phi$ to $\varsigma^{-1}(K)$ is proper.
\end{itemize}
\end{definition}

\begin{remark}\label{R-3-3}
For $m\geq 3$, one readily checks that the function
\[
\Phi(p)=\sum_{\substack{[v_i,v_j]\in E \\ v_k\in V\setminus\{v_i,v_j\}}}\frac{1}
{\left(\|p(v_i)-p(v_k)\|+\|p(v_j)-p(v_k)\|-\|p(v_i)-p(v_j)\|\right)^2}
\]
mentioned in Remark~\ref{R-1-10} serves as a $\varsigma$-potential on $\mathcal M(R_m)$.
\end{remark}




With these concepts, we acquire L-R functions on
$\mathcal M(R_m)$ as follows.

\begin{theorem}\label{T-3-5}
Let $m\geq 2$. Assume $\Phi:\mathcal M(R_m)\to [0,+\infty)$ is a $\varsigma$-potential and assume $h:\mathcal L(R_m)\to\mathbb R$ is a Lyapunov-Reeb function. Then
\[
f=\Phi-\Phi\circ\varsigma+h\circ\varsigma
\]
is a Lyapunov-Reeb function on $\mathcal M(R_m)$.
\end{theorem}

\begin{proof}
For each $p\in \mathcal M(R_m)$, since $\varsigma(p)$ is the straight configuration sharing the same length vector as $p$, we derive from
property~\textbf{(a1)} and Theorem~\ref{T-1-11} that
\begin{equation}\label{E-3-1}
\Phi(p)\geq \Phi\big(\varsigma(p)\big).
\end{equation}
Given that $h$ is lower-bounded,  $f$ must also be lower-bounded.

Next we prove $f$ is proper. For $c_1=\inf f$ and $c_2>\inf f$, it suffices to show that any sequence $\{p_n\}$ contains a convergent subsequence, provided
\begin{equation}\label{E-3-2}
c_1\leq f(p_n)=\Phi(p_n)-\Phi\big(\varsigma(p_n)\big)
+h\big(\varsigma(p_n)\big)\leq c_2.
\end{equation}
Because $\Phi$ satisfies property~\textbf{(a2)} and $h$ is a Lyapunov-Reeb function on $\mathcal L(R_m)$, we reduce the proof to verifying that both $\big\{h(\varsigma(p_n))\big\}$ and $\big\{\Phi(p_n)\big\}$ are upper-bounded. From~\eqref{E-3-1} and~\eqref{E-3-2}, it is straightforward to see
\[
h\big(\varsigma(p_n)\big)\leq c_2.
\]
By the properness of $h$, we deduce that $\big\{\varsigma(p_n)\big\}$ lies in the compact set
\[
K:=h^{-1}\big([b_2,c_2]\big)\subset \mathcal L(R_m),
\]
where $b_2=\min_{q\in \mathcal L(R_m)} h(q)$. Reapplying~\eqref{E-3-2}, we obtain
\[
\Phi(p_n)\leq c_2-b_2+\max\nolimits_{q\in K}\Phi(q).
\]

It remains to demonstrate that  $f$ possesses exactly one critical point. To do this, we regard $f$ as a function of $(\rho,\theta)$ and infer that
\[
f(\rho,\theta)=\Phi(\rho,\theta)-\Phi(\rho,\vec{0})+h(\rho,\vec{0}).
\]
Suppose $(\tilde{\rho},\tilde{\theta})$ is a critical point of $f$. For $i=2,\cdots,m-1$, taking the derivative with respect to $\theta_i$ gives
\[
\frac{\partial f}{\partial \theta_i}(\tilde{\rho},\tilde{\theta})=\frac{\partial \Phi}{\partial \theta_i}(\tilde{\rho},\tilde{\theta})=0,
\]
which implies $\tilde{\theta}$ is a critical point of  $\Phi(\tilde{\rho},\cdot)$. By property~\textbf{(a1)} and Theorem~\ref{T-1-11}, this yields $\tilde{\theta}=\vec{0}$. For $i=1,\cdots,m-1$, the derivative with respect to $\rho_i$ satisfies
\[
\frac{\partial f}{\partial \rho_i}(\tilde{\rho},\tilde{\theta})=\frac{\partial \Phi}{\partial \rho_i}(\tilde{\rho},\tilde{\theta})-\frac{\partial \Phi}{\partial \rho_i}(\tilde{\rho},\vec{0})+\frac{\partial h}{\partial \rho_i}(\tilde{\rho},\vec{0})=\frac{\partial h}{\partial \rho_i}(\tilde{\rho},\vec{0})=0.
\]
It follows that $\tilde{\rho}=\rho^\ast$, where $\rho^\ast$ is the unique critical point of $h(\cdot,\vec{0})$. Thus, $f$ has exactly one critical point $(\tilde{\rho},\tilde{\theta})=(\rho^\ast,\vec{0})$. This completes the proof.
\end{proof}

\begin{remark}\label{R-3-6}
As an example, we observe that
\[
h(q)=\sum\nolimits_{i=1}^{m-1}\log^2\|q(v_{i+1})-q(v_i)\|
\]
is a Lyapunov-Reeb function on $\mathcal L(R_m)$.
\end{remark}

We now turn to the topological characterizations of the relevant spaces.

\begin{proof}[\textbf{Proof of Theorem~\ref{T-1-3}}]
We assume $m\geq 3$, as the case $m=2$ is trivial. By virtue of Theorems~\ref{T-1-11}, ~\ref{T-3-5} and Remarks~\ref{R-1-10},~\ref{R-3-3},~\ref{R-3-6}, both
$\mathcal M(R_m,\ell)$ and $\mathcal M(R_m)$ admit Lyapunov-Reeb functions. In view of Theorem~\ref{T-1-7}, the conclusions follow.
\end{proof}

\begin{proof}[\textbf{Proof of Corollary~\ref{C-1-13}}]
It is a consequence of property $(iii)$ of Theorem~\ref{T-1-7}.
\end{proof}


\section{Moduli spaces of cycle linkages and configurations}\label{S-4}
We further investigate the moduli spaces of cycle linkages and configurations. As before, we need to equip
$\mathcal M^{+}(C_m)$ and $\mathcal M^{+}(C_m,\ell)$ with smooth structures and construct L-R functions. To ensure
$\mathcal M^{+}(C_m,\ell)\neq \emptyset$, we assume throughout this section that $\ell:E\to\mathbb R_{+}$ satisfies condition~$\mathrm{\mathbf{(c1)}}$.

\subsection{Smooth structures of moduli spaces}
We first note that $\mathcal M^{+}(C_m)$ is a non-empty open subset of
$\mathcal M(R_m)$, so it inherits the induced smooth structure from $\mathcal M(R_m)$.
Furthermore, we find that $\mathcal M^{+}(C_m,\ell)$ carries the smooth structure of an embedded submanifold of $\mathcal M^{+}(C_m)$.
For deeper results on this topic, the reader is referred to the seminal works of Kapovich and Millson~\cite{Kapovich-1995, Kapovich-1996,Kapovich-2002}.

\begin{lemma}\label{L-4-1}
Suppose $m\geq 3$ and suppose $\ell:E\to\mathbb R_{+}$ satisfies
condition~$\mathrm{\mathbf{(c1)}}$. The following hold:
\begin{itemize}
\item[$(i)$]$\mathcal M^{+}(C_m)$ is a smooth manifold of dimension $2m-3$
\item[$(ii)$]$\mathcal M^{+}(C_m,\ell)\subset \mathcal M^{+}(C_m)$ is a smooth submanifold of dimension $m-3$.
\end{itemize}
\end{lemma}

\begin{proof}
Since $\mathcal M^{+}(C_m)$ is a non-empty open subset of $\mathcal M(R_m)$, we deduce part $(i)$ from Lemma~\ref{L-3-1}. For part $(ii)$, we define a map $F:\mathcal M^{+}(C_m)\to \mathbb R_{+}^{m}$ by
\[
F(p)=\big(\|p(v_2)-p(v_1)\|,\cdots,\|p(v_m)-p(v_{m-1})\|,\|p(v_1)-p(v_m)\|\big).
\]
Then $\mathcal M^{+}(C_m,\ell)=F^{-1}(\vec{\ell})$, where
\[
\vec{\ell}=\big(\ell([v_1,v_2]),\cdots,\ell([v_{m-1},v_m]),\ell([v_{m},v_1])\big).
\]
Under condition $\mathrm{\mathbf{(c1)}}$, $\mathcal M^{+}(C_m,\ell)\neq \emptyset$. The lemma follows as a consequence of the Regular Value Theorem~\cite{Lee-2018}, provided that $\vec{\ell}$  is a regular value of the map $F$. It therefore suffices to show the Jacobian $DF(p)$ has full rank $m$ for each $p\in F^{-1}(\vec{\ell})$. By invoking the coordinates $(\rho,\theta)$ for $p$, we infer that $DF(p)$ has the form
\[
\left[\begin{array}{ccccccc}
1 & 0  & \cdots & 0 & 0 & \cdots & 0\\
0 & 1  & \cdots & 0 & 0 & \cdots & 0\\
& \vdots  & \ddots &   & & \vdots & \\
0 & 0  & \cdots & 1 & 0 & \cdots & 0 \\
\frac{\partial\rho_m}{\partial\rho_1} & \frac{\partial\rho_m}{\partial\rho_2}
& \cdots & \frac{\partial\rho_m}{\partial\rho_{m-1}} & \frac{\partial\rho_m}{\partial\theta_2}
& \cdots & \frac{\partial\rho_m}{\partial\theta_{m-1}}
\end{array}\right],
\]
where
\[
\begin{aligned}
&\rho_m:=\|p(v_1)-p(v_m)\|
=\sqrt{\left(\rho_1+\sum\nolimits_{j=2}^{m-1}
\rho_j\cos\theta_j\right)^2
+\left(\sum\nolimits_{j=2}^{m-1}\rho_j\sin\theta_j\right)^2}.
\end{aligned}
\]
We claim $DF(p)$ has rank $m$. Assume this does not hold. For $i=2,\cdots,m-1$, then
\[
\frac{\partial\rho_m}{\partial\theta_i}=
\frac{-\left(\rho_1+\sum\nolimits_{j=2}^{m-1}
\rho_j\cos\theta_j\right)\rho_i\sin\theta_i
+\left(\sum\nolimits_{j=2}^{m-1}\rho_j\sin\theta_j\right)
\rho_i\cos\theta_i}{\rho_m}=0,
\]
which is equivalent to saying the vector
\[
p(v_m)-p(v_1)
=\left(\rho_1+\sum\nolimits_{j=2}^{m-1}\rho_j\cos\theta_j,
\sum\nolimits_{j=2}^{m-1}\rho_j\sin\theta_j\right)
\]
is parallel to
\[
p(v_{i+1})-p(v_i)=\left(\rho_i\cos\theta_i,\rho_i\sin\theta_i\right).
\]
As a result, $p(v_1),p(v_2),\cdots,p(v_m)$ are collinear, contradicting that they are the vertices of a simple polygon.

We thus verify that $DF(p)$ has rank $m$, which concludes
 $\mathcal M^{+}(C_m,\ell)=F^{-1}(\vec{\ell})$ is an embedded smooth submanifold of $\mathcal M^{+}(C_m)$ of dimension $m-3$, as desired.
\end{proof}



\subsection{Polygon triangulations and local coordinates}
Next we introduce a system of local coordinates for the moduli spaces via polygon triangulations. To achieve this, we invoke the following classical result. For a proof, we refer the reader to the celebrated Two Ears Theorem attributed to Meisters~\cite{Meisters-1975}. See also the monograph by Devadoss and O'Rourke~\cite[Theorem 1.4]{Devadoss-2011} for an alternative approach.

\begin{lemma}\label{L-4-2}
Every simple polygon in $\mathbb R^2$ admits a triangulation without adding vertices. Moreover, if the polygon has $m$ vertices, then each of such triangulations has $m-2$ triangles and $m-3$ diagonals.
\end{lemma}


\begin{figure}[htbp]
     \centering
\begin{tikzpicture}[scale=0.7]
\coordinate (a) at (0.2,-1.3);
\coordinate (b) at (2.2,0);
\coordinate (c) at (2.4,2.2);
\coordinate (d) at (0.3,1.5);
\coordinate (e) at (-1.9,2.6);
\coordinate (f) at (-2.5,-0.2);
\coordinate (g) at (-1,0.18);
\coordinate (h) at (-1,1);
\coordinate (i) at (0,1.4);
\coordinate (k) at (0,3);
\draw (a)--(b)--(c)--(d)--(e)--(f)--(g)--(a);
\draw[dashed](d)--(b);
\draw[dashed](a)--(d);
\draw[dashed](d)--(g);
\draw[dashed](e)--(g);

\fill (a) circle (2pt);
\fill (b) circle (2pt);
\fill (c) circle (2pt);
\fill (d) circle (2pt);
\fill (e) circle (2pt);
\fill (f) circle (2pt);
\fill (g) circle (2pt);
\node[below] at (1.2,0.8) {$e$};
%\draw ([shift={(0.2,-1.1)}]0:0.4) arc[radius=0.5, start angle=45, end angle=90];
\draw ([shift={(0.21,-1)}]0:0.4)
arc[radius=0.5, start angle=32, end angle=80];
%\draw ([shift={(2.48,2.55)}]-125:0.6) arc[radius=0.3, start angle=-160, end angle= %-95];
\draw ([shift={(2.47,2.56)}]-125:0.6)
arc[radius=0.3, start angle=-166, end angle= -106];
\node[below] at (0.7,0) {$\gamma$};
\node[below] at (2,2) {$\beta$};

\end{tikzpicture}
\caption{Polygon triangulation with diagonals marked by dotted lines}
\label{F-2}
\end{figure}

We further present a refinement of Lemma~\ref{L-4-2},  serving as a tool for proving the area-increasing property in the next subsection (see Lemma~\ref{L-4-9}). Suppose
$\mathcal T$ is a triangulation of a simple polygon without adding vertices. A diagonal $e$ in $\mathcal T$ is called  a \textbf{Lawson edge} if $\beta+\gamma\leq \pi$, where $\beta,\gamma$ are the angles opposite $e$ in the two adjacent triangles of $\mathcal T$ (see Fig.~\ref{F-2}).

The lemma below follows from Lawson's Local Optimization Procedure~\cite{Lawson-1984}. Improved results can also be found in the works of De Floriani, Falcidieno and Pienovi~\cite{Floriani-1985}, Lee and Lin~\cite{Lee-1986}, Chew~\cite{Chew-1989}, Chin and Wang~\cite{Chin-1999}, and others.

\begin{lemma}\label{L-4-3}
Every simple polygon in $\mathbb R^2$ admits a triangulation such that no extra vertices are added and all diagonals are Lawson edges.
\end{lemma}

Recall that each $p\in \mathcal M^{+}(C_m)$ gives a simple polygon $p(C_m)$. By Lemma~\ref{L-4-2}, $p(C_m)$ admits a triangulation $\mathcal T$ without adding vertices. Let $U_{\mathcal T}\subset \mathcal M^{+}(C_m)$ be a neighborhood of $p$ such that every $p'\in U_{\mathcal T}$ yields a simple polygon admitting a triangulation of the same combinatorial type as $\mathcal T$. We thus obtain a map
\[
\begin{aligned}
\varphi_{\mathcal T}:\;&U_{\mathcal T}\quad \longrightarrow\qquad\qquad\quad \mathbb R_{+}^{2m-3}\\
&\;p\;\;\quad \longmapsto\;\left(l_1(p),\cdots,l_{m}(p),\lambda_{1}(p),\cdots,\lambda_{m-3}(p)\right),
\end{aligned}
\]
where $l_1(p),\cdots,l_{m}(p)$ denote the side lengths of $p(C_m)$, and $\lambda_{1}(p),\cdots,\lambda_{m-3}(p)$ denote the lengths of the diagonals in the triangulation. Since
$\varphi_{\mathcal T}$ is a diffeomorphism from $U_{\mathcal T}$ to an open subset $\varphi_{\mathcal T}(U_{\mathcal T}) \subset \mathbb R_{+}^{2m-3}$,
$(U_{\mathcal T}, \varphi_{\mathcal T})$ forms a coordinate chart on $\mathcal M^{+}(C_m)$.
Moreover, for another chart $(U_{\mathcal T'}, \varphi_{\mathcal T'})$ with
$U_{\mathcal T}\cap U_{\mathcal T'}\neq \emptyset$,  the transition map
\[
\varphi_{\mathcal T'} \circ\varphi^{-1}_{\mathcal T}~:~\varphi_{\mathcal T}\left(U_{\mathcal T}\cap U_{\mathcal T'}\right)
~\longrightarrow~\varphi_{\mathcal T'}\left(U_{\mathcal T}\cap U_{\mathcal T'}\right)
\]
is smooth. Let $\Lambda$ be the set of combinatorial types of triangulations of all positively oriented simple $m$-gons. Then
$\mathcal A=\left\{(U_{\mathcal T}, \varphi_{\mathcal T})\right\}_{\mathcal T\in \Lambda}$ constitutes a smooth atlas for $\mathcal M^{+}(C_m)$,
with $(l,\lambda):=(l_1,\cdots,l_{m}, \lambda_1,\cdots,\lambda_{m-3})$ as local coordinates. Note that
\[
\mathcal M^{+}(C_m,\ell)=\left\{(l,\lambda)\in\mathcal M^{+}(C_m):l=\vec{\ell}\right\}.
\]
This implies that $\lambda_{1},\cdots,\lambda_{m-3}$ serve as local coordinates for points in $\mathcal M^{+}(C_m,\ell)$.

\subsection{Properties of the area function}
We now explore several properties of the area function.
An element in $\mathcal M^{+}(C_m,\ell)$ (resp. $\mathcal M^{+}(C_m)$) is called a \textbf{cocircular linkage} (resp. \textbf{cocircular configuration}) if it corresponds to a simple polygon inscribed in a circle. By Penner's Theorem~\cite[Theorem 6.2]{Penner-1987}, there exists exactly one cocircular linkage in
$\mathcal M^{+}(C_m,\ell)$. Relevant results also appeared in the works of Schlenker~\cite{Schlenker-2007}, Kou\v{r}imsk\'{a}, Skuppin and Springborn~\cite{Kour-Sku-Springborn-2016}, Xu and Zhou~\cite{Xu-2024}, and many others.

The following proposition, rooted in classical geometry, traces back to early investigations by Moula~\cite{Moula-1737}, L'Huilier~\cite{Huilier-1782}, Steiner~\cite{Steiner-1841}, and  Blaschke~\cite{Blaschke-1956}. For completeness, we provide an independent proof using the local coordinates.

\begin{proposition}\label{P-4-4}
The area function $A:\mathcal M^{+}(C_m,\ell)\to\mathbb R_+$ has the cocircular linkage as its unique critical point.
\end{proposition}

We first consider some properties of the area function for triangles.
\begin{lemma}\label{L-4-5}
Suppose $\Delta$ is a triangle in $\mathbb R^2$ with side lengths $l_i,l_j,l_k$ and opposite angles $\alpha_i,\alpha_j,\alpha_k$, respectively. Let $A(\Delta)$ denote the area of $\Delta$. The following formulas hold:
\begin{itemize}
\item[$(i)$] $\displaystyle{\frac{\partial A(\Delta)}{\partial l_i}=\frac{l_i\cot\alpha_i}{2}}$;
\item[$(ii)$] $\displaystyle{\frac{\partial^2 A(\Delta)}{\partial l_i^2}=\frac{\cot\alpha_i}{2}
    -\frac{1}{2l_il_jl_k}\left(\frac{l_i}{\sin\alpha_i}\right)^3}$;
\item[$(iii)$] $\displaystyle{\frac{\partial^2 A(\Delta)}{\partial l_i \partial l_j}=\frac{1}{2l_il_jl_k}\left(\frac{l_i}{\sin\alpha_i}\right)^3\cos\alpha_k}$.
\end{itemize}
\end{lemma}

\begin{remark}\label{R-4-6}
Set $\displaystyle{\Upsilon_{ijk}=\frac{1}{2l_il_jl_k}
\left(\frac{l_i}{\sin\alpha_i}\right)^3}$.
By the Law of Sines,
\[
\Upsilon_{ijk}=\Upsilon_{jki}=\Upsilon_{kij}:=\Upsilon(\Delta).
\]
We present the following equivalent versions of the latter two formulas:
\begin{itemize}
\item[$(ii')$] $\displaystyle{\frac{\partial^2 A(\Delta)}{\partial l_i^2}=\frac{\cot\alpha_i}{2}-\Upsilon(\Delta)}$;
\item[$(iii')$] $\displaystyle{\frac{\partial^2 A(\Delta)}{\partial l_i \partial l_j}=\Upsilon(\Delta)\cos\alpha_k}$.
\end{itemize}
\end{remark}


The proof is based on direct computations, with details deferred to~\S\ref{S-7}. Using this lemma, we are ready to show Proposition~\ref{P-4-4}.


\begin{proof}[\textbf{Proof of Proposition~\ref{P-4-4}}]
We assume $m\geq 4$, as the case $m=3$ is trivial. For each $p\in \mathcal M^{+}(C_m,\ell)$, let $\mathcal T$ be a triangulation of the polygon $p(C_m)$ without adding vertices. Suppose $\lambda_1,\cdots,\lambda_{m-3}$ are the lengths of the diagonals in the triangulation.  For $i=1,\cdots,m-3$, it follows from Lemma~\ref{L-4-5} that
\begin{equation}\label{E-4-1}
\frac{\partial A}{\partial \lambda_i}=\frac{\lambda_i}{2}\left(\cot\beta_i+\cot\gamma_i\right)
=\frac{\lambda_i\sin(\beta_i+\gamma_i)}{2\sin \beta_i\sin\gamma_i},
\end{equation}
where $\beta_i, \gamma_i$ are the angles opposite the  $i$-th diagonal in the two adjacent triangles of $\mathcal T$. At the cocircular linkage, we have
\begin{equation}\label{E-4-2}
\beta_i+\gamma_i=\pi,
\end{equation}
which yields
\begin{equation}\label{E-4-3}
\frac{\partial A}{\partial \lambda_i}=0.
\end{equation}
Hence the cocircular linkage is a critical point of the area function.

Conversely, any critical point of $A$ must satisfy~\eqref{E-4-3}, which is equivalent to~\eqref{E-4-2} and  is fulfilled exclusively by the cocircular linkage. Since the cocircular linkage in $\mathcal M^{+}(C_m,\ell)$ is unique, the proposition follows.
\end{proof}


\begin{remark}%\label{R-4-7}
The area function actually attains its maximum at the cocircular linkage. This property can be derived using Steiner's idea of ``mirroring the dent''~\cite{Steiner-1838} combined with a compactness argument.
\end{remark}

The following proposition asserts the non-degeneracy of the critical point, which in turn enables us to establish the smoothness of the projection map in the next subsection (see Proposition~\ref{P-4-14}).

\begin{proposition}\label{P-4-8}
Suppose $m\geq 4$. The Hessian of $A:\mathcal M^{+}(C_m,\ell)\to \mathbb R_+$ at the critical point is negative definite.
\end{proposition}

\begin{proof}
By virtue of Proposition~\ref{P-4-4}, the only critical point of $A$ is the cocircular  linkage $q$, which yields a simple polygon $q(C_m)$ inscribed in a circle. This implies $q(C_m)$ is a strictly convex polygon, admitting a triangulation $\mathcal T$ with diagonals connecting vertices $v_{m-1}$ and $v_i$ for $i=1,\cdots,m-3$. As shown in Fig.~\ref{F-3}, we introduce the following notations:
\begin{itemize}
\item{} For $i=1,\cdots,m-3$, let $\lambda_i$ denote the length of the diagonal in $\mathcal T$ between vertices $v_{m-1}$ and $v_i$;
\item{} For $j=1,\cdots,m-2$, let $\Delta_j$ be the triangle in $\mathcal T$ with vertices $v_{j-1},v_{j},v_{m-1}$ (where $v_0:=v_m$), and let $\sigma_j,\delta_j,\phi_j$ represent the angles of  $\Delta_j$ at vertices $v_{j-1},v_j,v_{m-1}$, respectively.
\end{itemize}

\begin{figure}[htbp]
     \centering
\begin{tikzpicture}[scale=0.79]
\coordinate (b) at (1.6,-2.5);
\coordinate (c) at (3.36,-0);
\coordinate (d) at (2.9,2.3);
\coordinate (e) at (-0.2,3.9);
\coordinate (f) at (-2.81,2.41);
\coordinate (g) at (-3.35,-0.1);
\coordinate (a) at (-1.4,-2.6);

\draw (a)--(b)--(c)--(d)--(e)--(f)--(g)--(a)--(b);

\draw[dashed](f)--(b);
\draw[dashed](f)--(a);
\draw[dashed](f)--(d);
\draw[dashed](f)--(g);
\draw[dashed](f)--(c);

\fill (a) circle (2pt);
\fill (b) circle (2pt);
\fill (c) circle (2pt);
\fill (d) circle (2pt);
\fill (e) circle (2pt);
\fill (f) circle (2pt);
\fill (g) circle (2pt);
\node[below] at (-1.4,-2.7) {$v_1$};
\draw ([shift={(-1.4,-2.6)}]125:0.6) arc[radius=0.4, start angle=127, end angle= 98];
\draw ([shift={(-1.4,-2.6)}]104:0.45) arc[radius=0.4, start angle=100, end angle= 0];
\node[below] at (1.75,-2.6) {$v_2$};
\draw ([shift={(1.6,-2.5)}]180:0.4) arc[radius=0.4, start angle=170, end angle= 125];
\draw ([shift={(1.6,-2.5)}]130:0.4) arc[radius=0.4, start angle=130, end angle= 60];
\node[right] at (3.39,-0.1) {$v_3$};
\draw ([shift={(3.36,-0)}]235:0.4) arc[radius=0.4, start angle=230, end angle= 160];
\draw ([shift={(3.36,-0)}]160:0.5) arc[radius=0.4, start angle=170, end angle= 95];
\node[right] at (2.96,2.35) {$v_4$};
\draw ([shift={(2.9,2.3)}]-80:0.3) arc[radius=0.3, start angle=-80, end angle= -180];
\draw ([shift={(2.9,2.3)}]180:0.5) arc[radius=0.4, start angle=180, end angle= 150];
\node[above] at (-0.2,3.96) {$v_5$};
\draw ([shift={(-0.2,3.9)}]-30:0.3) arc[radius=0.3, start angle=-30, end angle= -150];
\node[left] at (-2.82,2.48) {$v_6$};
\draw ([shift={(-2.81,2.41)}]30:0.5) arc[radius=0.4, start angle=30, end angle= -5];
\draw ([shift={(-2.81,2.41)}]-2:0.7) arc[radius=0.4, start angle=-5, end angle=-35];
\draw ([shift={(-2.81,2.41)}]-23:0.5) arc[radius=0.4, start angle=-14, end angle=-44];
\draw ([shift={(-2.81,2.41)}]-50:0.65) arc[radius=0.6, start angle=-45, end angle=-70];
\draw ([shift={(-2.81,2.41)}]-75:0.5) arc[radius=0.6, start angle=-85, end angle=-105];
\node[left] at (-3.39,-0.1) {$v_7$};
\draw ([shift={(-3.35,-0.1)}]75:0.3) arc[radius=0.3, start angle=75, end angle=-52];
\node at (-2.74,1.25) {$\phi_1$};
\node at (-2.08,1.2) {$\phi_2$};
\node[below] at (-1.7,2.05) {$\phi_3$};
\node at (-1.3,2.12) {$\phi_4$};
\node at (-1.75,2.65) {$\phi_5$};

\node at (-2.4,-0.55) {$\Delta_1$};
\node at (-0.4,-1.4) {$\Delta_2$};
\node at (1.35,-0.45) {$\Delta_3$};
\node at (1.6,1.4) {$\Delta_4$};
\node at (0.5,3) {$\Delta_5$};

\node at (-1.88,-1.56) {$\delta_1$};
\node at (1,-2.2) {$\delta_2$};
\node at (2.7,-0.1) {$\delta_3$};
\node at (2.55,1.9) {$\delta_4$};
\node at (-0.2,3.35) {$\delta_5$};
\node at (-2.7,0.1) {$\sigma_1$};
\node at (-1.1,-1.95) {$\sigma_2$};
\node at (1.6,-1.9) {$\sigma_3$};
\node at (2.75,0.6) {$\sigma_4$};
\node at (1.8,2.55) {$\sigma_5$};
\node at (-1.75,-0.25) {$\lambda_1$};
\node at (-0.2,0) {$\lambda_2$};
\node at (0.45,1.45) {$\lambda_3$};
\node at (-0.3,2.6) {$\lambda_4$};


\draw[dashed] (0,0.5) circle (3.4);

\end{tikzpicture}
\caption{A cocircular heptagon with a triangulation}
\label{F-3}
\end{figure}

Recall that $\lambda_1,\cdots,\lambda_{m-3}$ serve as local coordinates. From Lemma~\ref{L-4-5} and Remark~\ref{R-4-6}, it follows that
\[
\frac{\partial^2 A}{\partial \lambda_i^2}=\frac{\partial^2 A(\Delta_i)}{\partial \lambda_i^2}+\frac{\partial^2 A(\Delta_{i+1})}{\partial \lambda_i^2}=\frac{\sin(\sigma_i+\delta_{i+1})}{2\sin\sigma_i\sin\delta_{i+1}}-\Upsilon(\Delta_i)-\Upsilon(\Delta_{i+1}).
\]
At the cocircular linkage $q$, we have $\sigma_i+\delta_{i+1}=\pi$, which yields
\[
\frac{\partial^2 A}{\partial \lambda_i^2}\big(q\big)=-\Upsilon(\Delta_i)-\Upsilon(\Delta_{i+1})<0.
\]
For $1\leq j<k\leq m-3$, Lemma~\ref{L-4-5} and Remark~\ref{R-4-6} also imply that
\[
\frac{\partial^2 A}{\partial \lambda_j\partial \lambda_k}\big(q\big)
=\begin{cases}
\begin{aligned}
&\Upsilon(\Delta_{j+1})\cos\phi_{j+1},&k-j=1,\\
&0,  &k-j>1.
\end{aligned}
\end{cases}
\]
Therefore,
\[
\left|\frac{\partial^2 A}{\partial \lambda_i^2}\big(q\big)\right|-\sum_{j\neq i}\left|\frac{\partial^2 A}{\partial \lambda_i\partial \lambda_j}\big(q\big)\right|>0.
\]
In summary, the Hessian of $A$ at $q$ is a strictly diagonally dominate matrix with negative diagonal entries, and hence negative definite.
\end{proof}


The following area-increasing property was established by Connelly, Demaine and Rote~\cite[Theorem 7]{Connelly-2003}. Here we provide a simpler proof based on Lemma~\ref{L-4-3}.

\begin{lemma}[Connelly-Demaine-Rote]\label{L-4-9}
Assume $\{p_t\}_{-\varepsilon<t<\varepsilon}$ $(\varepsilon>0)$ is an expansive motion in $\mathcal M^{+}(C_m,\ell)$. Then
\[
\frac{\mathrm{d}A(p_t)}{\mathrm{d}t}\geq0,\;\forall\, t\in (-\varepsilon,\varepsilon).
\]
\end{lemma}
\begin{proof}
By Lemma~\ref{L-4-3}, we obtain a triangulation $\mathcal T$ of $p_t(C_m)$ with no extra vertices added and all diagonals being Lawson edges. For $i=1,\cdots,m-3$, let $\lambda_{i}$ be the length of the $i$-th diagonal in $\mathcal T$, and let $\beta_i, \gamma_i$ be the angles opposite the $i$-th diagonal in the two adjacent triangles of
$\mathcal T$. Since all diagonals in $\mathcal T$ are Lawson edges, it follows that
$\beta_i+\gamma_i\leq \pi$. Combining this with~\eqref{E-4-1}, we derive
\[
\frac{\partial A}{\partial \lambda_i}=\frac{\lambda_i\sin(\beta_i+\gamma_i)}{2\sin \beta_i\sin \gamma_i}\geq 0.
\]
Recall that $\{p_t\}_{-\varepsilon<t<\varepsilon}$ is an expansive motion. Then
\[
\frac{\mathrm{d}\lambda_i}{\mathrm{d}t}\geq 0.
\]
A direct computation yields
\[
\frac{\mathrm{d}A(p_t)}{\mathrm{d}t}=\sum\nolimits_{i=1}^{m-3}\frac{\partial A}{\partial \lambda_i}\frac{\mathrm{d}\lambda_i}{\mathrm{d}t}\geq 0.
\]
\end{proof}

\subsection{Constructing Lyapunov-Reeb functions} Remember that each $p\in \mathcal M^{+}(C_m)$ corresponds to a simple polygon $p(C_m)$ in $\mathbb R^2$, with $\alpha_i(p)$ denoting the interior angle at vertex $p(v_i)$. As mentioned in~\S~\ref{S-1},  constructing L-R functions on $\mathcal M^{+}(C_m,\ell)$ requires analyzing the function
\[
w=\frac{1}{m}\left(\sum\nolimits_{i=1}^m w_i\right),
\]
where
\[
w_i(p)=
\begin{cases}
\begin{aligned}
&\exp\left(\frac{1}{\pi-\alpha_i(p)}\right), &&\alpha_i(p)>\pi,\\
&0, &&\alpha_i(p)\leq\pi.
\end{aligned}
\end{cases}
\]

Note that $w(p)=0$ if $p\in \mathcal M^{+}(C_m,\ell)$ is a \textbf{convex cycle linkage} (i.e., $p$ corresponds to a convex polygon). The following lemma is due to Shimamoto and Wootters~\cite{Shimamoto-2014}. For the sake of self-containment, we include a proof.


\begin{lemma}[Shimamoto-Wootters]\label{L-4-10}
Assume $\{p_t\}_{-\varepsilon<t<\varepsilon}$ $(\varepsilon>0)$ is an expansive motion in $\mathcal M^{+}(C_m,\ell)$. Then
\[
\frac{\mathrm{d}w(p_t)}{\mathrm{d}t}\bigg|_{t=0}\leq0.
\]
\end{lemma}

\begin{proof}
To establish the lemma, we proceed by showing that
\[
\frac{\mathrm{d}w_i(p_t)}{\mathrm{d}t}\bigg|_{t=0}\leq0
\]
for $i=1,\cdots,m$. If $\alpha_i(p_0)\leq \pi$, it is straightforward to see
\[
\frac{\mathrm{d}w_i(p_t)}{\mathrm{d}t}\bigg|_{t=0}=0.
\]
For $\alpha_i(p_0)>\pi$, the Law of Cosines implies
\[
\frac{\mathrm{d}\alpha_i(p_t)}{\mathrm{d}t}
=\frac{\|p_t(v_{i-1})-p_t(v_{i+1})\|}{l_il_{i+1}\sin \alpha_i(p_t)}\cdot\frac{\mathrm{d}}{\mathrm{d}t}~\|p_t(v_{i-1})-p_t(v_{i+1})\|,
\]
where $l_i=\ell\big([v_{i-1},v_{i}]\big)$, $l_{i+1}=\ell\big([v_{i},v_{i+1}]\big)$. Since $\{p_t\}_{-\varepsilon<t<\varepsilon}$ $(\varepsilon>0)$ is an expansive motion, the distance $\|p_t(v_{i-1})-p_t(v_{i+1})\|$ is non-decreasing, which yields
\[
\frac{\mathrm{d}\alpha_i(p_t)}{\mathrm{d}t}\bigg|_{t=0}\leq 0.
\]
As a result,
\[
\frac{\mathrm{d}w_i(p_t)}{\mathrm{d}t}\bigg|_{t=0}
=\exp\left(\frac{1}{\pi-\alpha_i(p_0)}\right)
\left(\frac{1}{\pi-\alpha_i(p_0)}\right)^2
\frac{\mathrm{d}\alpha_i(p_t)}{\mathrm{d}t}\bigg|_{t=0}\leq 0.
\]
This completes the proof.
\end{proof}


Furthermore, we invoke the following fundamental result due to Connelly, Demaine and Rote~\cite[Theorem 3]{Connelly-2003}. It is worth mentioning that an independent proof can be found in the work of Rote, Santos and Streinu~\cite[Theorem 4.3]{Rote-2003}.
\begin{theorem}[Connelly-Demaine-Rote]\label{T-4-11}
If $p\in \mathcal M^{+}(C_m,\ell)$ is not a convex cycle linkage, then there exists an expansive motion $\{p_t\}_{-\varepsilon<t<\varepsilon}$ ($\varepsilon>0$) in $\mathcal M^{+}(C_m,\ell)$ such that $p_0=p$.
\end{theorem}

Let us move on to showing Theorem~\ref{T-1-12}.

\begin{proof}[\textbf{Proof of Theorem~\ref{T-1-12}}]
First it is easy to see $f$ is lower-bounded. Next we prove the properness of $f$. For $c_1=\inf f\geq 0$ and $c_2>\inf f$, we need to verify that any sequence
\[
\{p_n\}\subset f^{-1}\big([c_1,c_2]\big)\subset \mathcal M^{+}(C_m,\ell)
\]
contains a convergent subsequence.
Because the map $\Phi:\mathcal M^{+}(C_m,\ell)\to [0,+\infty)$ is proper, it is equivalent to validating that $\big\{\Phi(p_n)\big\}$ has an upper bound.
Assume on the contrary that there exists a subsequence $\{p_{n_k}\}\subset \{p_n\}$ such that
\[
\Phi(p_{n_k})\to+\infty.
\]
This implies $\{p_{n_k}\}$ tends to a linkage with self-intersections. On the other hand, given that $\{p_{n_k}\}\subset f^{-1}([c_1,c_2])$, we have
\[
0\leq c_1\leq f(p_{n_k})=\frac{1}{A(p_{n_k})}+w(p_{n_k})\Phi(p_{n_k})\leq c_2.
\]
Consequently,
\[
A(p_{n_k})\geq \frac{1}{c_2}>0\quad\text{and}\quad w(p_{n_k})\to 0.
\]
Since $w(p)=0$ if and only if $\alpha_i(p)\leq \pi$ for $i=1,\cdots,m$, we derive that $\{p_{n_k}\}$ approaches a convex cycle linkage with positive area, which has no self-intersections. This leads to a contradiction, so $f$ must be proper.


To demonstrate that $f$ has a unique critical point, we decompose
$\mathcal M^{+}(C_m,\ell)$ into the union of two disjoint subsets:
$\mathcal M_1$ and $\mathcal M_2$, where
$\mathcal M_1$ consists of the convex cycle linkages in $\mathcal M^{+}(C_m,\ell)$, and $\mathcal M_2:=\mathcal M^{+}(C_m,\ell)\setminus\mathcal M_1$.

On $\mathcal M_1$, $w\equiv 0$, so the function $f$ reduces to $\displaystyle{1/A}$. By Proposition~\ref{P-4-4}, $f$ possesses a unique critical point on $\mathcal M_1$.

We claim, meanwhile, that $f$ has no critical points on $\mathcal M_2$. Assume, for contradiction, that $q\in \mathcal M_2$ is a critical point of $f$. Since $q$ is  non-convex, Theorem~\ref{T-4-11} guarantees the existence of an expansive motion $\{p_t\}_{-\varepsilon<t<\varepsilon}$ ($\varepsilon>0$) in $\mathcal M^{+}(C_m,\ell)$ with $p_0=q$. The derivative of $f$ along this motion is expressed as
\[
\frac{\mathrm{d}f(p_t)}{\mathrm{d}t}=
-\frac{1}{A^2}\frac{\mathrm{d}A(p_t)}{\mathrm{d}t}
+w(p_t)\frac{\mathrm{d}\Phi(p_t)}{\mathrm{d}t}+\Phi(p_t)\frac{\mathrm{d}w(p_t)}{\mathrm{d}t}.
\]
By Lemmas~\ref{L-4-9},~\ref{L-4-10} and definition of strain energies, we obtain
\[
\frac{\mathrm{d}f(p_t)}{\mathrm{d}t}\bigg|_{t=0}<0,
\]
which contradicts the assumption that $p_0=q$ is a critical point of $f$.

Therefore, $f$ has exactly one critical point on
$\mathcal M^{+}(C_m,\ell)$. Putting the arguments together, we prove the theorem.
\end{proof}


\begin{remark}\label{R-4-12}
From the above proof, we can see that the function $f=1/A+w\Phi$ attains its global minimum at the cocicular linkage.
\end{remark}



\begin{remark}\label{R-4-13}
Shimamoto and Wootters~\cite{Shimamoto-2014} introduced the function $F=w\Phi$ to study the moduli space. Their proof relied crucially on a compactness argument regarding the set of convex cycle linkages. However, compactness may fail if the constraint\textemdash that for every sign assignment $\sigma: E\to\{-1,1\}$ to the edges, $\sum_{e_i\in E}\sigma(e_i)\ell(e_i)\neq0$\textemdash is removed. Fortunately, by virtue of Theorem~\ref{T-1-7}, we need not resort to specific topological analysis of the set of convex cycle linkages.
\end{remark}


To proceed, we further construct L-R functions on the moduli space
$\mathcal M^{+}(C_m)$. For each $p\in \mathcal M^{+}(C_m)$, Penner's Theorem~\cite[Theorem 6.2]{Penner-1987} implies there exists a unique cocircular configuration
$\tau(p)\in \mathcal M^{+}(C_m)$ sharing the same length vector as $p$. This gives rise to the following projection map
\[
\begin{aligned}
\tau\;:\;\mathcal M^{+}(C_m)\;&\longrightarrow\;\mathcal Q^{+}(C_m)\\
p\;\;\;\,&\longmapsto\;\;\tau(p),
\end{aligned}
\]
where $\mathcal Q^{+}(C_m)\subset\mathcal M^{+}(C_m)$ denotes the subspace of cocircular configurations. It is necessary to analyze the smoothness of this map.

\begin{proposition}\label{P-4-14}
$\mathcal Q^{+}(C_m)\subset\mathcal M^{+}(C_m)$ is a smooth submanifold of dimension $m$. Moreover, the map $\tau:\mathcal M^{+}(C_m)\to\mathcal Q^{+}(C_m)$ is smooth.
\end{proposition}

\begin{proof}
Without loss of generality, we assume $m\geq 4$. Given $p\in\mathcal M^{+}(C_m)$, let $\mathcal T$ be a triangulation of the polygon $p(C_m)$ without adding vertices, and let $(l,\lambda)$ be the local coordinates associated to
$\mathcal T$. Note that $\tau(p)$ corresponds to a cocircular polygon admitting a triangulation combinatorially equivalent to
 $\mathcal T$. According to Proposition~\ref{P-4-4}, $\tau(p)$ is parameterized by $\big(l,\lambda^\ast(l)\big)$, where $\big(l,\lambda^\ast(l)\big)$ satisfies
\[
\frac{\partial A}{\partial \lambda_i}\big(l,\lambda^\ast(l)\big)=0\quad\text{for}\quad i=1,\cdots,m-3.
\]
By invoking Proposition~\ref{P-4-8} and the Implicit Function Theorem~\cite{Lee-2018}, we deduce that $\lambda^\ast(l)$ depends smoothly on $l$, which completes the proof.
\end{proof}

%As in the preceding section, we introduce the following %concept.

\begin{definition}
A smooth map $\Phi:\mathcal M^{+}(C_m)\to [0,+\infty)$ is called a \textbf{$\tau$-potential}
on $\mathcal M^{+}(C_m)$ if the following properties hold:
\begin{itemize}
\item[\textbf{(s1)}] For every $q\in \mathcal Q^{+}(C_m)$, the restriction of $\Phi$ to $\tau^{-1}(q)$ is a strain energy;
\item[\textbf{(s2)}]For every compact set $K\subset \mathcal Q^{+}(C_m)$, the restriction of $\Phi$ to $\tau^{-1}(K)$ is proper.
\end{itemize}
\end{definition}

We construct L-R functions on $\mathcal M^{+}(C_m)$ as follows.
\begin{theorem}\label{T-4-16}
Let $m\geq 3$. Assume $\Phi:\mathcal M^{+}(C_m)\to [0,+\infty)$ is a $\tau$-potential and assume $h:\mathcal Q^{+}(C_m)\to\mathbb R$ is a Lyapunov-Reeb function. Then
\[
f=\frac{1}{A}-\frac{1}{A\circ\tau}+w\Phi+h\circ\tau
\]
is a Lyapunov-Reeb function on $\mathcal M^{+}(C_m)$.
\end{theorem}
\begin{proof}
For every $p\in \mathcal M^{+}(C_m)$, $\tau(p)$ is automatically convex, and thus $w\big(\tau(p)\big)\equiv0$. In light of property~\textbf{(s1)}, Theorem~\ref{T-1-12} and Remark~\ref{R-4-12}, we have
\begin{equation}\label{E-4-4}
\frac{1}{A(p)}+w(p)\Phi(p)\geq \frac{1}{A\big(\tau(p)\big)}
+w\big(\tau(p)\big)\Phi\big(\tau(p)\big)
=\frac{1}{A\big(\tau(p)\big)}.
\end{equation}
Since $h$ is lower-bounded, $f$ must also be lower-bounded.

To verify that $f$ possesses exactly one critical point, we decompose
\[
\mathcal M^{+}(C_m)=\mathcal Z_1\cup \mathcal Z_2,
\]
where $\mathcal Z_1\subset \mathcal M^{+}(C_m)$ denotes the subset of convex cycle configurations, and
$\mathcal Z_2$ is defined as the complement $\mathcal M^{+}(C_m)\setminus\mathcal Z_1$.
On $\mathcal Z_1$, observe that
\[
f=\frac{1}{A}-\frac{1}{A\circ\tau}+h\circ\tau.
\]
Following the proof of Theorem~\ref{T-3-5},  we assert $f$ has a unique critical point on $\mathcal Z_1$. Meanwhile, analogous arguments to those in Theorem~\ref{T-1-12} imply that $f$ has no critical points on $\mathcal Z_2$. Therefore, $f$ has exactly one critical point on $\mathcal M^{+}(C_m)$.

It remains to show that $f$ is proper. For $c_1=\inf f$ and $c_2>c_1$, we need to demonstrate that any sequence $\{p_n\}$ satisfying
\begin{equation}\label{E-4-5}
c_1\leq f(p_{n})=\frac{1}{A(p_{n})}
-\frac{1}{A\big(\tau(p_n)\big)}+w(p_{n})\Phi(p_{n})
+h\big(\tau(p_n)\big)
\leq c_2
\end{equation}
contains a convergent subsequence. Combining~\eqref{E-4-4} and~\eqref{E-4-5} yields
\[
h\big(\tau(p_n)\big)\leq c_2.
\]
Given that $h$ is a Lyapunov-Reeb function, $\big\{\tau(p_n)\big\}$ lies within a compact subset of $\mathcal Q^{+}(C_m)$, so $\big\{A\big(\tau(p_n)\big)\big\}$ is bounded both below and above by positive constants. This implies that there exists $c_3>0$ such that
\begin{equation}\label{E-4-6}
0<\frac{1}{A(p_{n})}+w(p_{n})\Phi(p_{n})\leq c_3.
\end{equation}
In view of property~\textbf{(s2)}, to confirm $\{p_n\}$ contains a convergent subsequence, it suffices to check that $\big\{\Phi(p_{n})\big\}$ is upper-bounded. Due to~\eqref{E-4-6}, this follows by analogy with the proof of Theorem~\ref{T-1-12}. Hence $f$ is proper. In conclusion, we establish that $f$ is a Lyapunov-Reeb function on $\mathcal M^{+}(C_m)$.
\end{proof}

Note that $\mathcal Q^{+}(C_m)$ is diffeomorphic to a non-empty open convex set $\Omega\subset\mathbb R_{+}^m$,  where each vector $(l_1,\cdots,l_m)$ satisfies
$l_i<\sum\nolimits_{j\neq i} l_j$ for $i=1,\cdots,m$. Thus, theoretically,
$\mathcal Q^{+}(C_m)$ admits L-R functions. We now construct a specific L-R function that will be applied to developing algorithms for the Configuration Refolding Problem in the next section. For $i=1,\cdots,m$, the length coordinate $l_i$ serves as a smooth function on $\mathcal Q^{+}(C_m)$. We relate these to L-R functions on
$\mathcal Q^{+}(C_m)$ as follows.

\begin{lemma}\label{L-4-17}
Let $L=\sum_{i=1}^m l_i$. Then
\[
h=\log^2 L-\left(\sum\nolimits_{i=1}^{m}\log\sin\dfrac{2\pi l_i}{L}\right)
\]
is a Lyapunov-Reeb function on $\mathcal Q^{+}(C_m)$.
\end{lemma}
\begin{proof}
Clearly, $h$ is proper and lower-bounded. It remains to verify that $h$ has exactly one critical point. For $i=1,\cdots,m-1$, we introduce the variable substitution $\mu_i=l_i/L$, regard $h$ as a function of $(\mu,L):=(\mu_1,\cdots,\mu_{m-1},L)$, and derive
\[
h(\mu,L)=\log^2 L-\left(\sum\nolimits_{i=1}^{m-1}\log\sin2\pi\mu_i\right)
-\log\sin2\pi\left[1-\left(\sum\nolimits_{i=1}^{m-1}\mu_i\right)\right].
\]
A straightforward computation yields that $h$ has a unique critical point
\[
\big(\tilde{\mu},\tilde{L}\big):=\left(1/m,\cdots,1/m,1\right).
\]
Consequently, $h$ is a Lyapunov-Reeb function on $\mathcal Q^{+}(C_m)$.
\end{proof}

We are now in a position to prove Theorem~\ref{T-1-4} and Corollary~\ref{C-1-14}.

\begin{proof}[\textbf{Proof of Theorem~\ref{T-1-4}}]
We assume $m\geq 4$, as the case $m=3$ is trivial. Applying Theorems~\ref{T-1-7},~\ref{T-1-12},~\ref{T-4-16} and Lemma~\ref{L-4-17}, we reduce the proof to constructing a $\tau$-potential on $\mathcal M^{+}(C_m)$. As before, it is easy to verify that the function
\[
\Phi(p)=\sum_{\substack{[v_i,v_j]\in E \\ v_k\in V\setminus\{v_i,v_j\}}}\frac{1}{\left(\|p(v_i)-p(v_k)\|+\|p(v_j)-p(v_k)\|-\|p(v_i)-p(v_j)\|\right)^2}
\]
satisfies the required properties. This completes the proof.
\end{proof}

\begin{proof}[\textbf{Proof of Corollary~\ref{C-1-14}}]
The assertion follows from Theorem~\ref{T-1-7}.
\end{proof}


\section{Algorithms}\label{S-5}
This section is devoted to developing algorithms for the Carpenter's Rule Problem and the Linkage (Configuration) Refolding Problem. It is noteworthy that both problems hold significant importance in practical fields, particularly in computer graphics~\cite{Hopcroft-1985,Alt-2003} and robotic-arm folding mechanisms~\cite{Farber-2008}. Over recent decades, remarkable progress has been made in addressing these challenges. Seminal contributions include the works of Connelly, Demaine and Rote~\cite{Connelly-2003}, Streinu~\cite{Streinu-2005}, Cantarella, Demaine, Iben and O'Brien~\cite{Cantarella-2004}, Iben, O'Brien and Demaine~\cite{Demaine-2009}, and many others.

\subsection{Carpenter's Rule Problem}
First we study the problem of straightening arm linkages (configurations). Recall that $\mathcal M(R_m,\ell)$ is identified with an open subset $U_{m,\ell}\subset \left(\mathbb R/2\pi\mathbb Z\right)^{m-2}$ via the coordinates $\theta=(\theta_2,\cdots,\theta_{m-1})$. Let $m\geq 3$ and let $\Phi$ be given as in Remark~\ref{R-1-10}. Viewing $\Phi$ as a function of $\theta$, we infer from Theorem~\ref{T-1-11} that $\Phi$ is an L-R function on $U_{m,\ell}$ with $\vec{0}=(0,\cdots,0)\in U_{m,\ell}$ being the critical point. Consider the negative gradient flow
\begin{equation}\label{E-5-1}
\frac{\mathrm d\theta_i}{\mathrm d t}=-\frac{\partial \Phi}{\partial \theta_i},\quad i=2,\cdots,m-1.
\end{equation}

From Theorem~\ref{T-1-7}, we derive the following result.

\begin{theorem}
Suppose $m\geq 3$. The flow~\eqref{E-5-1} starting at any point in $U_{m,\ell}$ exists for all $t\geq 0$ and converges to $\vec{0}$ as $t\to+\infty$.
\end{theorem}

Let $\theta(t)$ be the solution to~\eqref{E-5-1} and let $p_t\in \mathcal M(R_m,\ell)$ be given by
\[
\begin{aligned}
& p_t(v_1)=(0,0),\quad p_t(v_2)=(\ell_1,0),\\
& p_t(v_i)=
\left(\ell_1+\sum\nolimits_{j=2}^{i-1}\ell_j\cos\theta_j(t),
\sum\nolimits_{j=2}^{i-1}\ell_j\sin\theta_j(t)\right),\quad i=3,\cdots,m,
\end{aligned}
\]
where $\ell_k=\ell\big([v_k,v_{k+1}]\big)$ for $k=1,\cdots,m-1$. Then $\{p_t\}_{0\leq t\leq +\infty}$ yields the desired straightening motion in $\mathcal M(R_m,\ell)$ (see Fig.~\ref{F-4}).

\begin{figure}[htbp]
    \centering
    \includegraphics[width=0.78\textwidth]{F4.jpg}
    \caption{Straightening an arm linkage}
    \label{F-4}
\end{figure}


Analogously, we identify $\mathcal M(R_m)$ with an open subset $U_m\subset \mathbb R_{+}^{m-1}\times\left(\mathbb R/2\pi\mathbb Z\right)^{m-2}$ via the coordinates $(\rho,\theta)=(\rho_1,\cdots,\rho_{m-1},\theta_2,\cdots,\theta_{m-1})$. Let $\Phi$ and $h$ be given as in Remarks~\ref{R-1-10} and~\ref{R-3-6}, respectively. By virtue of Theorem~\ref{T-3-5},
\[
f(\rho,\theta)=\Phi(\rho,\theta)-\Phi(\rho,\vec{0})+h(\rho,\vec{0})
\]
is an L-R function on $U_m$. Let $\big(\rho(t),\theta(t)\big)$ be the negative gradient flow of $f$ starting at a given vector in $U_m$, and let $p_t\in\mathcal M(R_m)$ be the configuration parameterized by $\big(\rho(t),\theta(t)\big)$. Then $\{p_t\}_{0\leq t\leq +\infty}$ forms a straightening motion in $\mathcal M(R_m)$.


We now address the problem of convexifying cycle linkages and configurations. To this end, let $\ell_i=\ell\big([v_i,v_{i+1}]\big)$ for $i=1,\cdots,m$ (with $v_{m+1}:=v_1$), and let
\[
\Xi=\left\{\theta=(\theta_2,\cdots,\theta_{m-1})\in (\mathbb R/2\pi\mathbb Z)^{m-2}:~u(\theta)=0\right\},
\]
where
\[
u(\theta):=\sqrt{\left(\ell_1
+\sum\nolimits_{j=2}^{m-1}\ell_j\cos\theta_j\right)^2
+\left(\sum\nolimits_{j=2}^{m-1}\ell_j\sin\theta_j\right)^2}-\ell_m.
\]
As in~\S\ref{S-4}, we identify $\mathcal M^{+}(C_m,\ell)$ with an open subset $W_{m,\ell}\subset\Xi$ and regard
\[
f=\frac{1}{A}+w\Phi
\]
as a function on $W_{m,\ell}$. Consider the flow
\begin{equation}\label{E-5-2}
\frac{\mathrm d\theta_i}{\mathrm d t}=-\frac{\partial f}{\partial \theta_i}+\frac{\langle \nabla f, \nabla u\rangle}{\|\nabla u\|^2}\frac{\partial u}{\partial \theta_i},
\quad i=2,\cdots,m,
\end{equation}
which projects the negative gradient of  $f$  onto the tangent space of the constraint manifold $\Xi$. Following the proof of Lemma~\ref{L-4-1}, we show that
$\|\nabla u\|^2\neq 0$ on $W_{m,\ell}$, ensuring the flow~\eqref{E-5-2} is well defined. Let $\theta(t)$ be the solution of~\eqref{E-5-2}. A routine calculation yields
\[
\frac{\mathrm d u\big(\theta(t)\big)}{\mathrm d t}=0,
\]
so $\theta(t)$ remains on $\Xi$ whenever
$\theta(0)\in W_{m,\ell}\subset \Xi$. Furthermore,
\[
\frac{\mathrm d f\big(\theta(t)\big)}{\mathrm d t}\leq 0,
\]
from which we conclude that $\theta(t)$ is confined to $W_{m,\ell}$ for all $t\geq 0$. Let $\theta^\ast$ be the critical point of $f$. Because of Theorem~\ref{T-1-7}, the following result holds.
\begin{theorem}\label{T-5-2}
Suppose $m\geq 3$ and suppose $\ell:E\to\mathbb R_{+}$ satisfies
condition~$\mathrm{\mathbf{(c1)}}$. The flow~\eqref{E-5-2} starting at any point in $W_{m,\ell}$ exists for all $t\geq 0$ and converges to $\theta^\ast$ as $t\to+\infty$.
\end{theorem}

Let $p_t\in \mathcal M^{+}(C_m,\ell)$ denote the linkage parameterized by $\theta(t)$. Thus $\{p_t\}_{0\leq t\leq +\infty}$ yields a convexifying motion in
$\mathcal M^{+}(C_m,\ell)$ (see Fig.~\ref{F-5}). Analogously, by invoking the L-R function described in Theorem~\ref{T-4-16}, one can derive a convexifying motion in $\mathcal M^{+}(C_m)$.

\begin{figure}[htbp]
    \centering
    \includegraphics[width=0.79\textwidth]{F5.jpg}
    \caption{Convexifying a cycle linkage}
    \label{F-5}
\end{figure}




\subsection{ Linkage (Configuration) Refolding Problem} In fact, constructing a motion between two given linkages (configurations) is closely related to the Carpenter's Rule Problem. For instance, one can deform a cycle linkage into a cocircular shape and then fold it back down to the shape of the second linkage. However, this method leads to a needless detour, particularly when the two linkages have ``near'' shapes. The Refolding Problem aims to find more direct motions between prescribed linkages (configurations).

We embed this problem within the framework of the following Renormalization Problem:
\emph{Let $(M,g)$ be a  Riemannian manifold of dimension $d\geq 1$, and let $f:M\to\mathbb R$ be a Lyapunov-Reeb function. Given $p_0,p_1\in M$, how can one construct a Riemannian metric $\tilde{g}$ on $M$ and a geodesic in $(M,\tilde{g})$ connecting  $p_0$ and $p_1$?} To address the problem, we define $\eta_{a,b}:\mathbb R\to\mathbb R$ as in Lemma~\ref{L-2-8}, where
\[
a=1+\max\big\{f(p_0),f(p_1)\big\}\quad \text{and}\quad b=2+\max\big\{f(p_0),f(p_1)\big\}.
\]
Let $\psi: \mathbb R\times M\to M$ denote the flow generated by the vector field $-\eta_{a,b}(f)\nabla f$. This means $\psi$ satisfies
\begin{equation}\label{E-5-3}
\frac{\mathrm d \psi(s,p)}{\mathrm d s}=-\eta_{a,b}\left(f\big(\psi(s,p)\big)\right)\nabla f\big(\psi(s,p)\big).
\end{equation}
Note that $-\eta_{a,b}(f)\nabla f$ is a compactly supported smooth vector field. By virtue of Theorem~\ref{T-2-1}, the flow $\psi$ is complete, so $\psi_s(\cdot):=\psi(s,\cdot):M\to M$ defines a diffeomorphism for every
$s\in \mathbb R$. Making use of the pullback metric $g_s:=\psi_s^\ast g$, we establish the following property.

\begin{proposition}\label{P-5-3}
There exists $s_0\geq0$ such that for any $s\geq s_0$, the given points $p_0,p_1$ can be connected by a geodesic in $(M,g_s)$. In particular, for any $\delta>0$, there is an integer $n_0\geq0$ such that for any $n\geq n_0$, the points $p_0,p_1$ can be connected by a geodesic in $(M,g_{n\delta})$.
\end{proposition}

\begin{proof}
Let $q\in M$ be the critical point of $f$, and let $B\subset M$ be a geodesic convex neighborhood of $q$ in $(M,g)$. Similar arguments to those in the proof of Theorem~\ref{T-2-5} show that there exists $s_0\geq0$ such that for any $s\geq s_0$,
\[
\psi_s(p_0)=\psi(s,p_0)\in B,\quad \psi_s(p_1)=\psi(s,p_1)\in B.
\]
Since $B\subset M$ is a geodesic convex set in $(M,g)$, one can find a geodesic arc $\gamma$ on $B$ joining the points $\psi_s(p_0)$ and $\psi_s(p_1)$. Thus, $p_0,p_1$ are connected by the curve
\[
\gamma_{s}:=\psi_{-s}(\gamma)=\psi(-s,\gamma),
\]
which is a geodesic arc in $(M,g_s)$. This completes the proof.
\end{proof}

With this proposition, we reduce the Renormalization Problem to constructing a geodesic $\gamma^0$ in $(M,g)$ for a pair of ``nearby'' points $\psi_{n_0\delta}(p_0)$ and $\psi_{n_0\delta}(p_1)$, followed by using the preimage curve $\psi_{-n_0\delta}(\gamma^0)$ (see Fig.~\ref{F-6}). For $i=0,1,2,\cdots$, define
\[
\gamma^i:=\psi_{-i\delta}(\gamma^0).
\]
A systematic approach to constructing the curve sequence $\{\gamma^0,\gamma^1,\cdots\}$ is achieved via the following iteration process:
\begin{equation}\label{E-5-4}
%\begin{cases}
%\,\gamma^0=\gamma,\\\,
\gamma^{i+1}=\psi_{-\delta}(\gamma^{i}).
%\end{cases}
\end{equation}
We hereby formalize the algorithm for the Renormalization Problem as follows:
\begin{itemize}
\item[$1.$] Select $\delta>0$ and determine the minimal non-negative integer $n_0$ such that $\psi_{n_0\delta}(p_0)$ and $\psi_{n_0\delta}(p_1)$ are connected by a geodesic $\gamma^0$ in $(M,g)$. In view of Proposition~\ref{P-5-3}, this can be achieved through a finite number of iterations and tests.
\item[$2.$] Implement the iterative process defined  in~\eqref{E-5-4} to derive the curve $\gamma^{n_0}$, which serves as a geodesic in $(M,g_{n_0\delta})$ connecting $p_0$ and $p_1$.
\item[$3.$] To improve precision, refine $\delta$ by choosing smaller values and repeat the preceding two steps.
\end{itemize}


\begin{remark}
The above algorithm relies crucially on the assumption that $f$ has exactly one critical point. Without this property, the first step may fail, as $\psi_{n\delta}(p_0)$ and $\psi_{n\delta}(p_1)$ might converge to distinct critical points. In fact, when compared with the method in~\cite{Demaine-2009}, the use of L-R functions exhibits the advantage of predicting the long-time behavior and streamlining the proof of convergence.
\end{remark}


\begin{figure}[htbp]
\centering
\begin{tikzpicture}[scale=0.67]
  \coordinate (inner left)  at (-0.7,-0.6);
  \coordinate (inner right) at (0.7,-0.6);
  \coordinate (outer left)  at (-2.9,4.8);
  \coordinate (outer right) at (2.5,5.2);
  \coordinate (q) at (0,-1.2);
  \draw[black] plot[smooth cycle,tension=0.5] coordinates {
    (-2.8,-1) (-4,5.5)  (-1,4.5) (0.3,4.5)
     (2,5.8) (3.8,5) (3.1,0.3)
    (2.0,-1.7) (-1,-2)
  };
  \draw[black] (0,-0.7) ellipse [x radius=1.0cm, y radius=1.0cm];
  \fill[black] (inner left) circle (2pt) node[left=5pt] {\small $\psi_{n_0\delta}(p_0)$};
  \fill[black] (inner right) circle (2pt) node[right=5pt] {\small $\psi_{n_0\delta}(p_1)$};
  \fill[black] (outer left) circle (2pt) node[left=2pt] {$p_0$};
  \fill[black] (outer right) circle (2pt) node[right=2pt] {$p_1$};
  \fill[black] (q) circle (2pt) node[right] {$q$};
  \draw[black] (inner left) .. controls (0,-0.5) ..
    node[pos=0.5, above=-0.9pt] {$\gamma^0$}
    (inner right);
  \draw[black] (outer left) .. controls (0,3.5) ..
    node[below=0.1pt] {$\gamma^{n_0}$}
    (outer right);
  \draw[dashed] (inner left) .. controls (-1.5,0.5) ..
    coordinate[pos=0.4] (leftDot1)
    coordinate[pos=0.6] (leftDot2)
    (outer left);
  \fill[black] (leftDot1) circle (2pt);
  \fill[black] (leftDot2) circle (2pt);
  \draw[dashed] (inner right) .. controls (1.5,0.5) ..
    coordinate[pos=0.4] (rightDot1)
    coordinate[pos=0.6] (rightDot2)
    (outer right);
  \fill[black] (rightDot1) circle (2pt);
  \fill[black] (rightDot2) circle (2pt);
  \draw[black] (leftDot1) .. controls (0,0.9) .. node[above] {$\gamma^1$}(rightDot1);
  \draw[black] (leftDot2) .. controls (0,1.7) ..
    node[above] (gamma2) {$\gamma^2$}
    (rightDot2);
  \fill[black] ([yshift=0.25cm]gamma2.north) circle (1pt);
  \fill[black] ([yshift=0.4cm]gamma2.north) circle (1pt);
  \fill[black] ([yshift=0.55cm]gamma2.north) circle (1pt);
\end{tikzpicture}
\caption{Renormalization process}\label{F-6}
\end{figure}


Let us consider the Arm Linkage Refolding Problem: Given $p_0,p_1\in \mathcal M(R_m,\ell)$, construct an ``optimal" motion between them. As before, we identify $\mathcal M(R_m,\ell)$ with an open subset $U_{m,\ell}\subset \left(\mathbb R/2\pi\mathbb Z\right)^{m-2}$ via the coordinates $\theta=(\theta_2,\cdots,\theta_{m-1})$ and equip $U_{m,\ell}$ with the Riemannian metric
\[
g=\sum\nolimits_{i=2}^{m-1}(\mathrm d\theta_i)^2.
\]
Let $\theta^0,\theta^1\in U_{m,\ell}$ be the coordinates of $p_0,p_1$, respectively, and let $f:U_{m,\ell}\to\mathbb R$ be given as in Theorem~\ref{T-1-11}. Suppose
$\psi:\mathbb R\times U_{m,\ell}\to U_{m,\ell}$ is constructed as in~\eqref{E-5-3}. For sufficiently small $\delta>0$, the map
$\psi_\delta: U_{m,\ell}\to U_{m,\ell}$ can be approximated by
\[
\begin{aligned}
D_\delta: U_{m,\ell}&\;\longrightarrow\; U_{m,\ell}\\
\theta\quad&\;\longmapsto\; \theta-\delta~\eta_{a,b}\big(f(\theta)\big)~\nabla f(\theta).
\end{aligned}
\]
In addition, $\psi_{-\delta}: U_{m,\ell}\to U_{m,\ell}$ is approximated by
\[
\begin{aligned}
D_{-\delta}: U_{m,\ell}&\;\longrightarrow\; U_{m,\ell}\\
\theta\quad&\;\longmapsto\; \theta+\delta~\eta_{a,b}\big(f(\theta)\big)~\nabla f(\theta).
\end{aligned}
\]
Utilizing the renormalization algorithm described above, we construct the desired motion between $p_0$ and $p_1$ (see Fig.~\ref{F-7}).


\begin{figure}[htbp]
    \centering
    \includegraphics[width=0.75\textwidth]{F7.jpg}
    \caption{Refolding arm linkages}
    \label{F-7}
\end{figure}


Similarly, we identify $\mathcal M(R_m)$ with a subset $U_m\subset \mathbb R_{+}^{m-1}\times\left(\mathbb R/2\pi\mathbb Z\right)^{m-2}$ and equip it with the Riemannian metric
\[
g=\sum\nolimits_{i=1}^{m-1}(\mathrm d\rho_i)^2+\sum\nolimits_{i=2}^{m-1}(\mathrm d\theta_i)^2.
\]
Leveraging the L-R function from Theorem~\ref{T-3-5}, we can derive the renormalization algorithm for the Arm Configuration Refolding Problem.


We proceed to address the problem of constructing an ``optimal'' motion between two points $p_0,p_1\in \mathcal M^{+}(C_m,\ell)$. Recall that $\mathcal M^{+}(C_m,\ell)$ is identified with a subset
\[
W_{m,\ell}\subset \Xi\subset (\mathbb R/2\pi\mathbb Z)^{m-2},
\]
where
\[
\Xi=\left\{\theta=(\theta_2,\cdots,\theta_{m-1})\in (\mathbb R/2\pi\mathbb Z)^{m-2}:~u(\theta)=0\right\},
\]
and
\[
u(\theta)=\sqrt{
\left(\ell_1+\sum\nolimits_{j=2}^{m-1}\ell_j\cos\theta_j\right)^2
+\left(\sum\nolimits_{j=2}^{m-1}\ell_j\sin\theta_j\right)^2}-\ell_m.
\]
The set $W_{m,\ell}$ is endowed with the induced metric from the Riemannian metric
\[
g=\sum\nolimits_{i=2}^{m-1}(\mathrm d\theta_i)^2
\]
of the ambient space $(\mathbb R/2\pi\mathbb Z)^{m-2}$. Let $\theta^0,\theta^1\in W_{m,\ell}$ be the coordinates of $p_0,p_1$, respectively, and let $f:W_{m,\ell}\to\mathbb R$ be an L-R function as in Theorem~\ref{T-1-12}. Consider the flow $\psi:\mathbb R\times W_{m,\ell}\rightarrow W_{m,\ell}$ generated by the projected vector field
\[
\eta_{a,b}(f)\left(-\nabla f+\frac{\langle \nabla f, \nabla u\rangle}{\|\nabla u\|^2}\nabla u\right).
\]
In light of Proposition~\ref{P-5-3}, for any $\delta>0$, there exists a non-negative integer $n_0$, such that $\psi_{n_0\delta}(\theta^0)$ and $\psi_{n_0\delta}(\theta^1)$ can be connected by a geodesic $\gamma^0$ in $W_{m,\ell}$ with respect to the induced metric. More precisely, the curve $\gamma^0$ is computable, as it satisfies the differential equation
\[
\frac{\mathrm d \gamma^0(t)}{\mathrm d t}=-\nabla H+\frac{\langle \nabla H, \nabla u\rangle}{\|\nabla u\|^2}\nabla u
\]
with the initial condition $\gamma^0(0)=\psi_{n_0\delta}(\theta^0)$, where
\[
H=\|\theta-\psi_{n_0\delta}(\theta^1)\|^2.
\]
Implementing the iteration process in~\eqref{E-5-4}, we obtain a geodesic
\[
\gamma^{n_0}:[0,+\infty]\to W_{m,\ell}
\]
that connects $\theta^0$ and $\theta^1$ under the renormalized metric. The family of cycle linkages parameterized by $\gamma^{n_0}$ then provides the desired motion between $p_0$ and $p_1$.

By analogy, we identify $\mathcal M^{+}(C_m)$ with a submanifold $W_m$ of $\mathbb R_{+}^m\times (\mathbb R/2\pi\mathbb Z)^{m-2}$ and construct the projected flow
$\psi: \mathbb R\times W_{m}\to W_{m}$ using the L-R function specified in Theorem~\ref{T-4-16}. Employing the renormalization algorithm, we thereby formulate a practical approach to the Cycle Configuration Refolding Problem (see Fig.~\ref{F-8}).

\begin{figure}[htbp]
    \centering
    \includegraphics[width=0.76\textwidth]{F8.jpg}
    \caption{Refolding cycle configurations}
    \label{F-8}
\end{figure}


\section{Questions}\label{S-6}
This paper leaves a number of questions open. We consider the following to be particularly interesting topics for further investigation.

1. \emph{Given $p\in \mathcal M^{+}(C_m)$, how to construct a Lyapunov-Reeb function on $\mathcal M^{+}(C_m)$ with $p$ as its critical point?} Analogous questions can also be posed for the moduli spaces $\mathcal M^{+}(C_m,\ell)$, $\mathcal M(R_m)$, and $\mathcal M(R_m,\ell)$. Note that the existence of such functions is trivial since these spaces are diffeomorphic to Euclidean spaces, but our goal is to derive explicit expressions for such functions. If this can be done, the associated negative gradient flows will yield more direct and efficient algorithms for the Linkage (Configuration) Refolding Problem.

2. \emph{Can one construct complete Riemannian metrics (with explicit expressions) on the moduli spaces $\mathcal M(R_m)$,  $\mathcal M(R_m,\ell)$, $\mathcal M^{+}(C_m)$, and $\mathcal M^{+}(C_m,\ell)$?} This problem is motivated by two objectives: $(1)$ Such metrics would facilitate the construction of geodesics between any pair of points in the moduli spaces, thereby providing a thorough resolution to the Linkage (Configuration) Refolding Problem; $(2)$ The associated volume forms and Brownian motions may offer insights into the path integral formulation in quantum physics (see~\cite{Charles-2010}). For related research on random polygons, refer to the works of Pardon~\cite{Pardon-2011,Pardon-2012}.


3. Let $\mathcal M^{+}(C_m,\mathbb H^2)$ denote the set of isometry classes of positively oriented, simple $m$-gons in the hyperbolic plane $\mathbb H^2$. \emph{What can be said about the topology of  $\mathcal M^{+}(C_m,\mathbb H^2)$?} Analogous questions arise for the moduli spaces
$\mathcal M(R_m,\mathbb H^2)$, $\mathcal M(R_m,\ell, \mathbb H^2)$,
and $\mathcal M^{+}(C_m,\ell, \mathbb H^2)$. A pivotal step in addressing these problems is to establish the local expansive property for hyperbolic linkages, which requires deriving results parallel to Theorems~\ref{T-3-2} and~\ref{T-4-11}. This endeavor proves challenging, however, because the linear programming method\textemdash particularly the dual principle\textemdash resists generalization to hyperbolic geometry.


4. Let $\mathcal M^{+}(C_\infty)$ represent the set of isometry classes of positively oriented, rectifiable simple closed curves in $\mathbb R^2$. \emph{Is it possible to construct a Lyapunov-Reeb function on $\mathcal M^{+}(C_\infty)$?} In view of Pardon's work~\cite{Pardon-2009}, $\mathcal M^{+}(C_\infty)$ is path-connected. Yet there is a need to uncover deeper topological properties of $\mathcal M^{+}(C_\infty)$ and devise practical algorithms for deforming rectifiable simple closed curves in $\mathbb R^2$. A natural approach to this end is constructing L-R functions on $\mathcal M^{+}(C_\infty)$.


\section{Appendix}\label{S-7}

In this section, we provide the proof of Lemma~\ref{L-4-5}.

\begin{proof}[\textbf{Proof of Lemma~\ref{L-4-5}}]
Consider the expression for the area
\[
A(\Delta)=\frac{1}{2}l_jl_k\sin \alpha_i.
\]
Taking the derivative with respect to $l_i$ yields
\begin{equation}\label{E-7-1}
\frac{\partial A(\Delta)}{\partial l_i}=\frac{1}{2}l_jl_k\cos \alpha_i\frac{\partial \alpha_i}{\partial l_i}.
\end{equation}
Meanwhile, from the Law of Cosines
\[
\cos \alpha_i=\frac{l_j^2+l_k^2-l_i^2}{2l_jl_k},
\]
it follows that
\begin{equation}\label{E-7-2}
\frac{\partial \alpha_i}{\partial l_i}=\frac{l_i}{l_jl_k\sin \alpha_i}.
\end{equation}
Combining~\eqref{E-7-1} and~\eqref{E-7-2}, we have
\begin{equation}\label{E-7-3}
\frac{\partial A(\Delta)}{\partial l_i}=\frac{l_i\cot\alpha_i}{2},
\end{equation}
which concludes part $(i)$.

For part $(ii)$, differentiating both sides of~\eqref{E-7-3} with respect to $l_i$ gives
\[
\frac{\partial A^2(\Delta)}{\partial l_i^2}=\frac{\cot\alpha_i}{2}-\frac{l_i}{2\sin^2 \alpha_i}\frac{\partial \alpha_i}{\partial l_i}
\]
Substituting~\eqref{E-7-2} into the above expression, we derive
\[
\frac{\partial^2 A(\Delta)}{\partial l_i^2}=\frac{\cot\alpha_i}{2}
    -\frac{1}{2l_il_jl_k}\left(\frac{l_i}{\sin\alpha_i}\right)^3.
\]

It remains to consider part $(iii)$. By virtue of~\eqref{E-7-3}, we infer
\[
\frac{\partial^2 A(\Delta)}{\partial l_i\partial l_j}=-\frac{l_i}{2\sin^2 \alpha_i}\frac{\partial \alpha_i}{\partial l_j}.
\]
As in previous steps, the Law of Cosines yields
\[
\frac{\partial \alpha_i}{\partial l_j}=-\frac{l_i\cos\alpha_k}{l_jl_k\sin \alpha_i}.
\]
Combining the above two formulas, we obtain
\[
\frac{\partial A^2(\Delta)}{\partial l_i\partial l_j}=\frac{1}{2l_il_jl_k}\left(\frac{l_i}{\sin \alpha_i}\right)^3\cos\alpha_k.
\]
This completes the proof.
\end{proof}

\section*{Acknowledgments}
Te Ba was supported by NSFC Grant No. 11631010, Ze Zhou was supported by NSFC Grant No. 12422105 and Grant No. 12371075.


\begin{thebibliography}{50}

\bibitem{Aichholzer-2003}
{\sc O.~Aichholzer, D.~Bremner, E.~D. Demaine, H.~Meijer, V.~Sacrist\'an, and
  M.~Soss}, {\em Long proteins with unique optimal foldings in the {H}-{P}
  model}, Comput. Geom., 25 (2003), pp.~139--159.

\bibitem{Aichholzer-2001}
{\sc O.~Aichholzer, E.~D. Demaine, J.~Erickson, F.~Hurtado, M.~Overmars,
  M.~Soss, and G.~T. Toussaint}, {\em Reconfiguring convex polygons}, Comput.
  Geom., 20 (2001), pp.~85--95.

\bibitem{Alt-2003}
{\sc H.~Alt, C.~Knauer, G.~Rote, and S.~Whitesides}, {\em The complexity of
  (un)folding}, in Proc. 19th Annu. Symp. Comput. Geom. (SCG 2003), ACM, 2003,
  pp.~164--170.

\bibitem{Barba-1954}
{\sc E.~A. Barba\v{s}in and N.~N. Krasovski\u{i}}, {\em On stability of motion
  in the large}, Doklady Akad. Nauk SSSR (N.S.), 86 (1952), pp.~453--456.

\bibitem{Blaschke-1956}
{\sc W.~Blaschke}, {\em Kreis und {K}ugel}, Walter de Gruyter \& Co., Berlin,
  1956.
\newblock 2nd edition.

\bibitem{Brown-1961}
{\sc M.~Brown}, {\em The monotone union of open {$n$}-cells is an open
  {$n$}-cell}, Proc. Amer. Math. Soc., 12 (1961), pp.~812--814.

\bibitem{Cantarella-2004}
{\sc J.~H. Cantarella, E.~D. Demaine, H.~N. Iben, and J.~F. O'Brien}, {\em An
  energy-driven approach to linkage unfolding}, in Proceedings of the 20th
  Annual Symposium on Computational Geometry, 2004.

\bibitem{Charles-2010}
{\sc L.~Charles}, {\em On the quantization of polygon spaces}, Asian J. Math.,
  14 (2010), pp.~109--152.

\bibitem{Chew-1989}
{\sc L.~P. Chew}, {\em Constrained Delaunay triangulations}, vol.~4, 1989,
  pp.~97--108.
\newblock Computational geometry (Waterloo, ON, 1987).

\bibitem{Chin-1999}
{\sc F.~Chin and C.~A. Wang}, {\em Finding the constrained {D}elaunay
  triangulation and constrained {V}oronoi diagram of a simple polygon in linear
  time}, SIAM J. Comput., 28 (1999), pp.~471--486.

\bibitem{Connelly-2009}
{\sc R.~Connelly and E.~D. Demaine}, {\em Geometry and topology of polygonal
  linkages}, in Handbook of discrete and computational geometry, Discrete
  Mathematics and its Applications (Boca Raton), CRC Press, third~ed., 2018,
  pp.~233--256.

\bibitem{Connelly-2003}
{\sc R.~Connelly, E.~D. Demaine, and G.~Rote}, {\em Straightening polygonal
  arcs and convexifying polygonal cycles}, Discrete Comput. Geom., 30 (2003),
  pp.~205--239.

\bibitem{Floriani-1985}
{\sc L.~{De Floriani}, B.~Falcidieno, and C.~Pienovi}, {\em Delaunay-based
  representation of surfaces defined over arbitrarily shaped domains}, Computer
  Vision, Graphics, and Image Processing, 32 (1985), pp.~127--140.

\bibitem{Demaine-2007}
{\sc E.~D. Demaine and J.~O'Rourke}, {\em Geometric folding algorithms.
  {Linkages}, origami, polyhedra}, Cambridge: Cambridge University Press, 2007.

\bibitem{Devadoss-2011}
{\sc S.~L. Devadoss and J.~O'Rourke}, {\em Discrete and computational
  geometry}, Princeton University Press, Princeton, NJ, 2011.


\bibitem{Farber-2008}
{\sc M.~Farber}, {\em Invitation to topological robotics}, Zurich Lectures in
  Advanced Mathematics, European Mathematical Society (EMS), Z\"urich, 2008.

\bibitem{Farber-Fromm-2012}
{\sc M.~Farber and V.~Fromm}, {\em The topology of spaces of polygons}, Trans.
  Amer. Math. Soc., 365 (2013), pp.~3097--3114.

\bibitem{Farber-Hausmann-2011}
{\sc M.~Farber, J.-C. Hausmann, and D.~Sch\"utz}, {\em The {W}alker conjecture
  for chains in {$\mathbb{R}^d$}}, Math. Proc. Cambridge Philos. Soc., 151
  (2011), pp.~283--292.

\bibitem{Farber-2007}
{\sc M.~Farber and D.~Sch\"utz}, {\em Homology of planar polygon spaces}, Geom.
  Dedicata, 125 (2007), pp.~75--92.


\bibitem{Fortier Bourque-2024}
{\sc M.~Fortier Bourque}, {\em The space of immersed polygons},
\newblock Preprint, {arXiv}:2406.02519.%2303.15406, 2023.arXiv:


\bibitem{Frank-Kamenetskii-1997}
{\sc M.~D. Frank-Kamenetskii}, {\em Unravelling {DNA}}, Addison-Wesley, 1997.

\bibitem{Freedman-1982}
{\sc M.~H. Freedman}, {\em The topology of four-dimensional manifolds}, J.
  Differential Geometry, 17 (1982), pp.~357--453.

\bibitem{Gonzalez-2016}
{\sc A.~Gonz\'{a}lez and J. L\'{o}pez-L\'{o}pez}, {\em Compactness of spaces of convex and simple quadrilaterals},
Bull. Aust. Math. Soc. 94 (2016), pp.~507--521.



\bibitem{Gonzalez-2018}
{\sc A.~Gonz\'{a}lez}, {\em Fundamental group of spaces of simple polygons},
  Rocky Mt. J. Math., 48 (2018), pp.~1871--1886.



\bibitem{Gonzalez-2023}
{\sc A.~Gonz{\'a}lez and M.~Sedano-Mendoza}, {\em Moduli space of simple
  polygons},
\newblock Preprint, {arXiv}:2303.15406, 2023.

\bibitem{Han-2013}
{\sc L.~Han, L.~Rudolph, M.~Chou, S.~Corbett, E.~Eagle, D.~Glotzer, J.~Kramer,
  J.~Moran, C.~Pietras, A.~Tareen, and M.~Valko}, {\em Configurations and path
  planning of convex planar polygonal loops}, in Algorithmic Foundations of
  Robotics X, E.~Frazzoli, T.~Lozano-Perez, N.~Roy, and D.~Rus, eds., Springer
  Berlin Heidelberg, 2013, pp.~53--69.

\bibitem{Hausmann-1989}
{\sc J.-C. Hausmann}, {\em Sur la topologie des bras articul\'es}, in Algebraic
  topology {P}ozna\'n{} 1989, vol.~1474 of Lecture Notes in Math., Springer,
  Berlin, 1991, pp.~146--159.

\bibitem{Hopcroft-1985}
{\sc J.~Hopcroft, D.~Joseph, and S.~Whitesides}, {\em On the movement of robot
  arms in {$2$}-dimensional bounded regions}, SIAM J. Comput., 14 (1985),
  pp.~315--333.

\bibitem{Demaine-2009}
{\sc H.~N. Iben, J.~F. O'Brien, and E.~D. Demaine}, {\em Refolding planar
  polygons}, Discrete Comput. Geom., 41 (2009), pp.~444--460.

\bibitem{Kamiyama-1998}
{\sc Y.~Kamiyama}, {\em The homology of singular polygon spaces}, Canad. J.
  Math., 50 (1998), pp.~581--594.

\bibitem{Kapovich-1995}
{\sc M.~Kapovich and J.~Millson}, {\em On the moduli space of polygons in the
  {E}uclidean plane}, J. Differential Geom., 42 (1995), pp.~430--464.

\bibitem{Kapovich-1996}
{\sc M.~Kapovich and J.~Millson}, {\em The symplectic geometry of polygons
  in {E}uclidean space}, J. Differential Geom., 44 (1996), pp.~479--513.

\bibitem{Kapovich-2002}
{\sc M.~Kapovich and J.~Millson}, {\em Universality theorems for
  configuration spaces of planar linkages}, Topology, 41 (2002),
  pp.~1051--1107.

\bibitem{Khimshiashvili-2009}
{\sc G.~Khimshiashvili}, {\em Configuration spaces and signature formulas},
  vol.~160, 2009, pp.~727--736.
\newblock Algebra and geometry.

\bibitem{Kirby-1997}
{\sc R.~Kirby}, {\em Problems in low-dimensional topology}, in Geometric
  topology: Proceedings of the 1993 Georgia International Topology Conference,
  W.~H. Kazez, ed., vol.~2 of AMS/IP Studies in Advanced Mathematics, American
  Mathematical Society, 1997, pp.~35--473.

\bibitem{Kour-Sku-Springborn-2016}
{\sc H.~Kou\v{r}imsk\'a, L.~Skuppin, and B.~Springborn}, {\em A variational
  principle for cyclic polygons with prescribed edge lengths}, in Advances in
  discrete differential geometry, Springer, Berlin, 2016, pp.~177--195.

\bibitem{Lawson-1984}
{\sc C.~Lawson}, {\em Software for {$C1$} surface interpolation}, in
  Mathematical Software, J.~R. Rice, ed., Academic Press, 1977, pp.~161--194.

\bibitem{Lee-1986}
{\sc D.~T. Lee and A.~K. Lin}, {\em Generalized {D}elaunay triangulation for
  planar graphs}, Discrete Comput. Geom., 1 (1986), pp.~201--217.

\bibitem{Lee-2018}
{\sc J.~Lee}, {\em Introduction to smooth manifolds}, vol.~218 of Graduate
  Texts in Mathematics, Springer, New York, second~ed., 2013.

\bibitem{Lenhart-1995}
{\sc W.~J. Lenhart and S.~H. Whitesides}, {\em Reconfiguring closed polygonal
  chains in {E}uclidean {$d$}-space}, Discrete Comput. Geom., 13 (1995),
  pp.~123--140.

\bibitem{Huilier-1782}
{\sc S.~L'Huilier}, {\em De relatione mutua capacitatis et terminorum
  figurarum, geometrice considerata: seu de maximis et minimis}, Sumptibus et
  typis Michaelis Gr{\"o}ll, 1782.

\bibitem{Lumelsky-1987}
{\sc V.~Lumelsky}, {\em Effect of kinematics on motion planning for planar
  robot arms moving amidst unknown obstacles}, IEEE Journal on Robotics and
  Automation, 3 (1987), pp.~207--223.

\bibitem{Meisters-1975}
{\sc G.~H. Meisters}, {\em Polygons have ears}, Amer. Math. Monthly, 82 (1975),
  pp.~648--651.

\bibitem{Milgram-2004}
{\sc R.~J. Milgram and J.~C. Trinkle}, {\em The geometry of configuration
  spaces for closed chains in two and three dimensions}, Homology Homotopy
  Appl., 6 (2004), pp.~237--267.

\bibitem{Milnor-1964}
{\sc J.~Milnor}, {\em Differential topology}, in Lectures on {M}odern
  {M}athematics, {V}ol. {II}, Wiley, New York-London-Sydney, 1964,
  pp.~165--183.

\bibitem{Moula-1737}
{\sc F.~Moula}, {\em De maximis in figuris rectilineis}, in Commentarii
  Academiae Scientiarum Imperialis Petropolitanae, vol.~9, Typis Academiae,
  1744, pp.~138--159.

\bibitem{Newman-1966}
{\sc M.~H.~A. Newman}, {\em The engulfing theorem for topological manifolds},
  Ann. of Math. (2), 84 (1966), pp.~555--571.

\bibitem{Pardon-2009}
{\sc J.~Pardon}, {\em On the unfolding of simple closed curves}, Trans. Amer.
  Math. Soc., 361 (2009), pp.~1749--1764.

\bibitem{Pardon-2011}
%\leavevmode\vrule height 2pt depth -1.6pt width 23pt,
{\sc J.~Pardon},{\em Central limit theorems for random polygons in an arbitrary convex set}, Ann. Probab. 39 (2011), pp.~881--903.

\bibitem{Pardon-2012}
%\leavevmode\vrule height 2pt depth -1.6pt width 23pt,
{\sc J.~Pardon},
{\em Central limit theorems for uniform model random polygons},
J. Theoret. Probab. 25 (2012), pp.~823--833.


\bibitem{Penner-1987}
{\sc R.~C. Penner}, {\em The decorated {T}eichm\"uller space of punctured
  surfaces}, Comm. Math. Phys., 113 (1987), pp.~299--339.

\bibitem{Perelman-2002}
{\sc G.~Perelman}, {\em The entropy formula for the {Ricci} flow and its
  geometric applications},
\newblock Preprint, {arXiv}:0211159, 2002.

\bibitem{Perelman-2003.3}
%\leavevmode\vrule height 2pt depth -1.6pt width 23pt,
{\sc G.~Perelman},{\em Finite extinction
  time for the solutions to the {Ricci} flow on certain three-manifolds},
\newblock Preprint, {arXiv}:0307245, 2003.

\bibitem{Perelman-2003.7}
%\leavevmode\vrule height 2pt depth -1.6pt width 23pt,
{\sc G.~Perelman}, {\em Ricci flow with
  surgery on three-manifolds},
\newblock Preprint, {arXiv}:0303109, 2003.

\bibitem{Reeb-1946}
{\sc G.~Reeb}, {\em Sur les points singuliers d'une forme de {P}faff
  compl\`etement int\'egrable ou d'une fonction num\'erique}, C. R. Acad. Sci.
  Paris, 222 (1946), pp.~847--849.

\bibitem{Rote-2003}
{\sc G.~Rote, F.~Santos, and I.~Streinu}, {\em Expansive motions and the
  polytope of pointed pseudo-triangulations}, in Discrete and computational
  geometry, The Goodman-Pollack Festschrift, Berlin: Springer, 2003,
  pp.~699--736.

\bibitem{Schlenker-2007}
{\sc J.-M. Schlenker}, {\em Small deformations of polygons and polyhedra},
  Trans. Amer. Math. Soc., 359 (2007), pp.~2155--2189.

\bibitem{Schuetz-2010}
{\sc D.~Sch\"utz}, {\em The isomorphism problem for planar polygon spaces}, J.
  Topol., 3 (2010), pp.~713--742.

\bibitem{Schuetz-2013}
%\leavevmode\vrule height 2pt depth -1.6pt width 23pt,
{\sc D.~Sch\"utz}, {\em Homology of moduli spaces of linkages in high-dimensional Euclidean space},
 Algebr. Geom. Topol. 13 (2013), pp.~1183--1224.

\bibitem{Schuetz-2016}
%\leavevmode\vrule height 2pt depth -1.6pt width 23pt,
{\sc D.~Sch\"utz},{\em Intersection homology of linkage spaces},
 J. Topol. Anal. 8 (2016), pp.~25-58.


\bibitem{Shimamoto-2005}
{\sc D.~Shimamoto and C.~Vanderwaart}, {\em Spaces of polygons in the plane and
  {M}orse theory}, Amer. Math. Monthly, 112 (2005), pp.~289--310.

\bibitem{Shimamoto-2014}
{\sc D.~Shimamoto and M.~Wootters}, {\em Configuration spaces of convex and
  embedded polygons in the plane}, Geom. Dedicata, 172 (2014), pp.~121--134.

\bibitem{Smale-1961}
{\sc S.~Smale}, {\em Generalized {P}oincar\'e's conjecture in dimensions
  greater than four}, Ann. of Math. (2), 74 (1961), pp.~391--406.

\bibitem{Stallings-1960}
{\sc J.~Stallings}, {\em Polyhedral homotopy-spheres}, Bull. Amer. Math. Soc.,
  66 (1960), pp.~485--488.

\bibitem{Stallings-1962}
{\sc J.~Stallings}, {\em The piecewise-linear
  structure of {E}uclidean space}, Proc. Cambridge Philos. Soc., 58 (1962),
  pp.~481--488.

\bibitem{Steiner-1838}
{\sc J.~Steiner}, {\em Einfache {B}eweise der isoperimetrischen
  {H}aupts\"atze}, J. Reine Angew. Math., 18 (1838), pp.~281--296.

\bibitem{Steiner-1841}
%\leavevmode\vrule height 2pt depth -1.6pt width 23pt,
{\sc J.~Steiner}, {\em Sur le maximum et le
  minimum des figures dans le plan, sur la sph{\`e}re, et dans l'espace en
  g{\'e}n{\'e}ral.}, J. Reine Angew. Math., 24 (1842), pp.~190--250.

\bibitem{Streinu-2005}
{\sc I.~Streinu}, {\em Pseudo-triangulations, rigidity and motion planning},
  Discrete Comput. Geom., 34 (2005), pp.~587--635.

\bibitem{Wilson-1967}
{\sc F.~Wilson}, {\em The structure of the level surfaces of a {L}yapunov
  function}, J. Differential Equations, 3 (1967), pp.~323--329.

\bibitem{Xu-2024}
{\sc Y.~Xu and Z.~Zhou}, {\em Polygons inscribed in {J}ordan curves with
  prescribed edge ratios}, Topology Appl., 354 (2024), pp.~Paper No. 108971, 9.

\bibitem{Zeeman-1961}
{\sc E.~Zeeman}, {\em The {P}oincar\'e{} conjecture for {$n\geq 5$}}, in
  Topology of 3-manifolds and related topics ({P}roc. {T}he {U}niv. of
  {G}eorgia {I}nstitute, 1961), Prentice-Hall, Inc., Englewood Cliffs, NJ,
  1961, pp.~198--204.

\end{thebibliography}


\end{document}
